% In-place centered-normalized proof; information-theoretic simplification, 2026-10-23.
% Template: main/template.tex; Chinese adaptation retains Tufte/Fandol fixes.
\documentclass[a4paper]{tufte-handout}
\usepackage{amsmath,amssymb,amsthm,booktabs,longtable}
\usepackage{fontspec,xeCJK}
\setmainfont{texgyrepagella-regular.otf}[BoldFont=texgyrepagella-bold.otf,ItalicFont=texgyrepagella-italic.otf,BoldItalicFont=texgyrepagella-bolditalic.otf]
\setCJKmainfont{FandolSong-Regular.otf}[BoldFont=FandolSong-Bold.otf,ItalicFont=FandolKai-Regular.otf]
\setCJKsansfont{FandolHei-Regular.otf}
\setCJKmonofont{FandolFang-Regular.otf}
\renewcommand{\allcaps}[1]{#1}
\renewcommand{\smallcaps}[1]{#1}
\renewcommand{\textsc}[1]{#1}
\setcounter{secnumdepth}{2}
\setlength{\emergencystretch}{3em}
\let\newthoughtwithbreak\newthought
\RenewDocumentCommand{\newthought}{s m}{\IfBooleanTF{#1}{\noindent\textbf{#2}}{\newthoughtwithbreak{#2}}}
\newtheorem{theorem}{定理}
\newtheorem{lemma}{引理}
\newtheorem{proposition}{命题}
\newcommand{\E}{\mathbb E}
\newcommand{\R}{\mathbb R}
\newcommand{\I}{\mathcal I}
\newcommand{\Pall}{P_{\mathrm{hidden}}}
\newcommand{\Pout}{P_{\mathrm{out}}}
\newcommand{\dsm}{\delta_{\mathrm{sm}}}
\newcommand{\ecut}{\eta_{\mathrm{cut}}}
\newcommand{\norm}[1]{\left\|#1\right\|}
\DeclareMathOperator{\tr}{tr}
\DeclareMathOperator{\KL}{KL}
\DeclareMathOperator{\Sym}{Sym}
\title{\texorpdfstring{中心化归一化网络的预测 KL 标度\\固定深度标度定理}{中心化归一化网络的预测 KL 标度：固定深度标度定理}}
\author{Gaussian 回归；维度先行极限}
\date{2026 年 10 月 23 日}
\hypersetup{unicode=true,hypertexnames=false,pdftitle={中心化归一化网络的预测 KL 标度：固定深度标度定理}}
\begin{document}
\maketitle
\begin{abstract}
固定深度、各层宽度与输入维数成正比时，中心化归一化 tanh、sigmoid、ReLU 的平均 Bayes 预测 KL 具有首项 $\Pall(d)/(2n)$。证明不假设固定样本比下的风险极限存在，也不假设后验局部 Gaussian 性。激活的高阶 Hermite 分量给出全局函数比较与有限标签参数匹配。辅助 Gaussian 序列通道将匹配转为下界所需的弱谱结论；条件 I--MMSE 与 Gaussian 分部积分给出预测下界，一致的线性 Gaussian 参考分布与辅助信息付费给出相同首项的上界。技术附录展开局部估计与跨网络函数比较。
\end{abstract}
\begin{fullwidth}

\section{模型、激活类与定理}
\newthought*{中心化和归一化是模型的明确假设。}
固定 $L\ge1$，$k_0=d$，$k_\ell/d\to\gamma_\ell\in(0,\infty)$，$\gamma_0=1$。
第 $\ell$ 层权重为 $W^{(\ell)}\in\R^{k_\ell\times k_{\ell-1}}$，其元素独立标准 Gaussian。输入 $x\sim N(0,I_d)$ 独立于参数。本文每层使用同一固定激活 $\sigma_\ell\equiv\sigma$，且
\[
\E\sigma(G)=0,\qquad\E\sigma(G)^2=1,\qquad G\sim N(0,1).
\]
记 $W=(W^{(\ell)})_{\ell=1}^L$、$\theta=\theta(W)=(v,W)$，固定 $v_i=1$，网络定义为
\begin{equation}\begin{aligned}
h^{(0)}(x)&=x,\qquad
h^{(\ell)}(x)=\sigma_\ell\!\left(\frac{W^{(\ell)}h^{(\ell-1)}(x)}{\sqrt{k_{\ell-1}}}\right),\\
F_\theta^{(L)}(x)&=\frac{v^{\mathsf T}h^{(L)}(x)}{\sqrt{k_L}}.
\end{aligned}\end{equation}
真实参数 $\theta^0=(v,W^0)$ 由上述先验 $P_0$ 抽取。训练集 $\mathcal D_n=\{(x_\mu,y_\mu)\}_{\mu=1}^n$ 由
\[
y_\mu=F_{\theta^0}^{(L)}(x_\mu)+\sqrt\Delta z_\mu,\qquad
z_\mu\sim N(0,1),\quad\Delta>0
\]
生成。训练和测试输入独立同分布；噪声相互独立，且独立于参数与全部输入。记 Gaussian 通道为
\[
\Pout(y\mid\lambda)=N(y;\lambda,\Delta).
\]
学习者使用真实先验和通道，其后验为
\[
P(d\theta\mid\mathcal D_n)\propto P_0(d\theta)
\prod_{\mu=1}^n\Pout\!\left(y_\mu\mid F_\theta^{(L)}(x_\mu)\right).
\]
只有隐藏矩阵被学习，读出 $v$ 固定；隐藏参数数目为
\[
\Pall(d)=\sum_{\ell=1}^Lk_{\ell-1}k_\ell,\qquad
c=\lim_d\frac{\Pall(d)}{d^2}
=\sum_{\ell=1}^L\gamma_{\ell-1}\gamma_\ell.
\]
后验预测密度及平均预测 KL 定义为
\begin{equation}\begin{aligned}
p_{\mathcal D_n}(y\mid x)
&=\E_{\theta\mid\mathcal D_n}\Pout\!\left(y\mid F_\theta^{(L)}(x)\right),\\
\epsilon_{\log,d,n}
&=\E_{\theta^0,\mathcal D_n,x}\KL\!\left(
\Pout\!\left(\cdot\mid F_{\theta^0}^{(L)}(x)\right)
\,\middle\|\,p_{\mathcal D_n}(\cdot\mid x)\right).
\end{aligned}\end{equation}
写 $\epsilon_{\log,d}(\alpha):=\epsilon_{\log,d,\lfloor\alpha d^2\rfloor}$。
技术证明中简记 $p=\Pall(d)$，并把 $W$ 视为 $p$ 维隐藏权重坐标；梯度、Hessian、Gaussian 积分和 OU 耦合均只作用于这些坐标。写 $\theta'=\theta(V)$ 表示第二组权重对应的参数。先验平均中可将真实隐藏权重 $W^0$ 简记为 $W$；技术附录中的 $F_W=F_{\theta(W)}^{(L)}$ 与 $F=F_W$ 均指同一个实际输出，除非另行明确指定齐次比较器。

所有常数允许依赖固定深度、激活、宽度比例与噪声，不能依赖 $d,n$。

\subsection{激活类}
\newthought*{取下列两类的并集 $\mathcal A_{\rm bal}$。}
两支均满足上述中心化及单位方差条件。
\emph{光滑奇函数支：} $\sigma$ 有界、奇、属于 $C^\infty$，每个固定阶导数有界。
\emph{分段线性支：} 令 $c_0=\sqrt{2/\pi}$，
\begin{equation}\label{eq:pl}
\sigma(z)=az+b(|z|-c_0),\qquad
b\ne0,\qquad a^2+b^2(1-2/\pi)=1.
\end{equation}
两支均全局 Lipschitz。条件只涉及一个实函数，不包含 Fisher、后验或信息极限前提。

记 $\operatorname{He}_j$ 为概率论约定的 Hermite 多项式，$\psi_j=\operatorname{He}_j/\sqrt{j!}$、$a_j=\E[\sigma(G)\psi_j(G)]$。
与网络符号一致，$\mu_{\ell,j}=\mu_j=\sqrt{j!}a_j$。以下事实使光滑支无需另加高阶系数假设。
\begin{lemma}[高阶奇系数自动存在]\label{lem:automatic-hermite}
若 $\sigma$ 为有界连续奇函数且 $\E\sigma(G)^2=1$，则存在奇数 $m\ge5$ 使 $a_m\ne0$。
\end{lemma}
\begin{proof}
奇性使所有偶次系数为零。若所有奇数 $m\ge5$ 的系数也为零，
Hermite 系的 $L^2(\gamma)$ 完备性给出
$\sigma=a_1\psi_1+a_3\psi_3$，Gaussian 几乎处处成立。
两边连续且 Gaussian 密度处处为正，故处处相等。
有界多项式只能为常数，奇性又使该常数为零，与单位方差矛盾。
\end{proof}
光滑支固定选择上述一个 $m$；分段线性支取 $m=6$，其 $a_6\ne0$ 在下节直接验证。
$m$ 可依赖固定激活，但不随维数或样本数变化；后文常数允许依赖此选择。

\begin{theorem}\label{thm:main}
上述模型中，任一固定 $\sigma\in\mathcal A_{\rm bal}$ 满足
\begin{equation}\label{eq:main}
\boxed{\lim_{\alpha\to\infty}\limsup_{d\to\infty}
\left|\alpha \epsilon_{\log,d}(\alpha)-\frac c2\right|=0.}
\end{equation}
该结论包含任意固定正斜率的中心化归一化 tanh、logistic sigmoid，以及中心化归一化 ReLU。
\end{theorem}
先取 $d\to\infty$，再取 $\alpha\to\infty$。不要求固定 $\alpha$ 的风险极限存在；不据此扩展到未中心化模型、层间任意混用两支、增长深度或非 Gaussian 输出的尖锐上界。

\section{三个激活为何属于此类}
\newthought*{只需验证上述标量函数性质。}
对任意固定 $a>0$，令
\[
s_a^2=\E\tanh^2(aG)>0,\qquad \sigma_a(z)=\tanh(az)/s_a.
\]
它有界、光滑且为奇函数，因而中心化，定义给出单位方差。
由 $\tanh'=1-\tanh^2$ 递推，每个固定阶导数都是 $\tanh(az)$ 的多项式乘固定常数，故有界。
因此 $\sigma_a$ 属于光滑奇函数支，高阶非零系数由引理~\ref{lem:automatic-hermite} 自动保证。

对任意固定 $b>0$，logistic sigmoid $\ell_b(z)=(1+e^{-bz})^{-1}$ 满足
\[
\ell_b(z)-\tfrac12=\tfrac12\tanh(bz/2).
\]
其均值为 $1/2$，方差为 $s_{b/2}^2/4$，故中心化归一化后恰为 $\sigma_{b/2}$。

对 ReLU，直接积分给出
\[
\mu_R=\E G_+=\frac1{\sqrt{2\pi}},\qquad
v_R=\operatorname{Var}(G_+)=\frac12-\frac1{2\pi},\qquad
\sigma_R(z)=\frac{z_+-\mu_R}{\sqrt{v_R}}.
\]
由 $z_+=(z+|z|)/2$，它对应式~\eqref{eq:pl} 的
$a=b=1/(2\sqrt{v_R})$。
整个分段线性支都中心化；奇偶正交性给出其方差
$a^2+b^2(1-2/\pi)=1$。
又由
\[
\operatorname{He}_6(z)=z^6-15z^4+45z^2-15,\qquad
\E|G|^{2j+1}=2^j j!\sqrt{2/\pi},
\]
得到 $\E[|G|\operatorname{He}_6(G)]=3\sqrt{2/\pi}$，所以
\[
a_6=\frac{3b\sqrt{2/\pi}}{\sqrt{6!}}\ne0.
\]
分段线性支的奇次系数 $a_j$ 在 $j\ge3$ 时均为零；光滑奇函数支的偶次系数均为零。
这正是后文两类剖面所使用的奇偶结构。

\section{先验矩与局部估计}
\newthought*{先建立不使用学习结论的基本控制。}
令 $q_\ell=\norm{h^{(\ell)}}^2/k_\ell$，$q_0=\norm X^2/d$。对固定 $r>0$，所需结论为
\begin{equation}\label{eq:moments}\begin{aligned}
\E|q_\ell-1|^r&\le C_r d^{-r/2},\qquad
\sup_d\E|F_\theta^{(L)}(X)|^r<\infty,\\
\sup_d\E\norm{\nabla_WF_\theta^{(L)}(X)}^r&<\infty.
\end{aligned}\end{equation}

输出与范数矩由附录引理 \ref{an:moments} 证明；其中保留了全一读出的条件均值
$\sqrt{k_L}\,\mathbb E[\sigma(\sqrt{q_{L-1}}G)]$。
梯度矩及光滑支的 $\mathbb E\|\nabla_W^2F\|_{\rm HS}^2\le C$
由引理 \ref{an:derivatives} 的前向、反向递推与 Hessian 分解给出。

两支统一取参数 OU 平滑
$F_{\dsm}(W,X)=\mathbb E_ZF(e^{-\dsm^2}W+\sqrt{1-e^{-2\dsm^2}}Z,X)$。
引理 \ref{an:smoothing} 证明
\begin{equation}\label{eq:smoothing}\begin{aligned}
\E\norm{\nabla F_{\dsm}-\nabla F}^2&\le C\dsm,\\
\E\norm{\nabla^2F_{\dsm}}_{\rm HS}^2&\le C/\dsm,\qquad
\E\norm{\nabla F_{\dsm}}^4\le C.
\end{aligned}\end{equation}
该平滑仅构造统计测试场，真实网络与似然不变。

记 $\theta(t)=\theta(W(t))$。令 $W(t)=\sqrt{1-t}W+\sqrt tZ$ 为保持先验的 OU 耦合。Gaussian Hermite 展开给出
\begin{equation}\label{eq:ou-rem}\begin{aligned}
F_{\theta(t)}^{(L)}-F_\theta^{(L)}&=\sqrt t\,(\nabla_W F_\theta^{(L)})^TZ+E_t,\\
\E\norm{E_t}_{L_X^2}^2&\le Ct^{1+\delta},\qquad
\delta=\begin{cases}1,&\text{光滑支},\\1/2,&\text{分段线性支}.
\end{cases}
\end{aligned}\end{equation}
这是完整的平稳余项，包含 OU 径向漂移与边界贡献的抵消。它仅用于特意构造的先验 OU 对，不用于任意后验对。

\section{全局函数比较与参数匹配}
\newthought*{同一匹配结论同时服务于信息上下界。}
写 $\psi_j=\operatorname{He}_j/\sqrt{j!}$、$a_j=\E\sigma(G)\psi_j(G)$，并置
\[
N=m^L,\qquad S=1+m+\cdots+m^{L-1},\qquad B=a_m^S\ne0.
\]
以各层齐次激活 $z^m/\sqrt{m!}$、原权重与原读出构造多项式网络，其最高输入次数系数为 $C_N(W)\in\Sym^N\R^d$。定义
\[
H_m(W)=I_N(C_N(W)),\qquad T_W=BH_m(W).
\]
这里 $I_N$ 是 Gaussian 多重积分，满足
$\langle I_N(C),I_N(D)\rangle_{L_X^2}=N!\langle C,D\rangle$。
所有张量使用真正的全局对称 Hilbert 范数；不会把分组张量范数当作实际函数范数。

\begin{lemma}[实际函数的耦合一致比较]\label{lem:cross}
若 $\Pi_d$ 为原 Gaussian 参数先验，则
\begin{equation}\label{eq:cross}
e_d:=\sup_{\pi\in\Gamma(\Pi_d,\Pi_d)}
\E_\pi|\langle T_W,F_{\theta'}^{(L)}-T_V\rangle_{L_X^2}|\longrightarrow0.
\end{equation}
这里允许任意正确先验边缘的耦合，不要求 $W,V$ 独立。
\end{lemma}
\begin{proof}[证明结构与实际接合]
先在各自先验下逼近真实的 $N$ 阶输入混沌。令
\[
P_N(z)=\sum_{j=1}^N\frac{a_jz^j}{\sqrt{j!}},\quad
Q_N(z)=\sum_{j=1}^N\frac{\beta_jz^j}{\sqrt{j!}},
\quad\beta_j=\E[(\sigma(G)^2-1)\psi_j(G)].
\]
令 $R_{t,N}[A;P_N]$ 为原深度 $t$ 前缀、根多项式 $A$、下层多项式 $P_N$ 的\emph{形式 $N$ 次系数}经 $I_N$ 实现的函数。它不是未截断多项式网络的整个 $N$ 阶投影。
光滑奇支有
\[
\Pi_NF_\theta^{(L)}=R_{L,N}[P_N;P_N]+o_{L^2(W,X)}(1).
\]
分段线性支令 $\mu=bc_0/2$、$U_N=[Q_N-P_N^2]_{\le N}$，则实际公式为
\begin{equation}\label{eq:rawclosure}\begin{aligned}
\Pi_NF_\theta^{(L)}={}&R_{L,N}[P_N;P_N]\\
&+\mu\sum_{t=1}^{L-1}\sqrt{k_L/k_t}\,R_{t,N}[U_N;P_N]
+o_{L^2(W,X)}(1).
\end{aligned}\end{equation}
$U_N$ 可含奇次项，不能强加普通分段线性节点的奇偶性。式 \eqref{eq:rawclosure} 的集体均值和二次恢复在附录 \ref{app:closure} 推导。

固定次数和深度后只剩有限个有序次数树及物理槽位分配。两支各自的最低切割归纳证明：与平衡教师配对时，深度 $L$ 唯一保留全 $m$ 次树；每个深度 $t<L$ 的任意正根次数源均为 $o(1)$。每个节点贡献 $a_m$，共有 $S$ 个节点，故保留项恰为 $T_V$。切割归纳及分段线性二次例外见附录 \ref{app:cuts}。

有限树估计在两个独立定义的边缘好事件上确定性成立；坏事件由固定高阶矩控制，因而对所有耦合一致。若实际逼近误差为 $E_d(V)$，则
\[
\E_\pi|\langle T_W,E_d(V)\rangle|
\le(\E\norm{T_W}^2)^{1/2}(\E\norm{E_d(V)}^2)^{1/2}=o(1).
\]
这只用边缘，不用两网络独立性。因 $T_W$ 位于第 $N$ 阶混沌，所有其余输入次数正交，得到 \eqref{eq:cross}。
\end{proof}

\subsection{从函数比较到参数匹配}
\newthought*{先极化，再匹配。}
对某个正确边缘耦合令
$D=\E\norm{T_W-T_V}^2$、$B_F=\E\norm{F_\theta^{(L)}-F_{\theta'}^{(L)}}^2$。
在原耦合、转置耦合以及两个对角项中分别使用 \eqref{eq:cross}，得
\[
|\E\langle T_W-T_V,F_\theta^{(L)}-F_{\theta'}^{(L)}\rangle-D|\le4e_d.
\]
所以
\[
D\le\sqrt{DB_F}+4e_d\le(D+B_F)/2+4e_d,
\quad D\le B_F+8e_d.
\]
齐次全局匹配定理给出有限可测映射族 $\mathcal G_d$，
\begin{equation}\label{eq:matching}\begin{aligned}
\log|\mathcal G_d|&=O(d\log d)=o(p),\\
\frac1p\E\min_{g\in\mathcal G_d}\norm{W-gV}^2
&\le C\E\norm{H_m(W)-H_m(V)}^2+o_d(1)\\
&\le C\E\norm{F_\theta^{(L)}-F_{\theta'}^{(L)}}^2+o_d(1).
\end{aligned}\end{equation}
误差对耦合一致。附录 \ref{app:matching} 说明如何逐层匹配而不损失误差指数。映射可含行符号；它们不必是真实激活的对称性。后文只把 $gV$ 当作 Gaussian 混合中心，绝不使用 $F_{\theta(gV)}^{(L)}=F_{\theta'}^{(L)}$。

\section{实际谱到预测 KL 下界}
\newthought*{Gaussian 通道无需额外的通道截断。}
令 $a_W(x)=\nabla_WF_\theta^{(L)}(x)$，
\[
A_W=\E_Xa_Wa_W^T,\qquad\I_W=A_W/\Delta,
\qquad D_d(r)=\E\frac1p\tr\{r(p\I_W+rI)^{-1}\}.
\]
\subsection{全局匹配的信息论推论}
令 $\mathcal H=L^2(P_X)$，$f(W)=F_\theta^{(L)}$，$J_W=D_Wf$，
$A_W=J_W^*J_W$。已有解析估计保证 $f\in W^{1,2}(\gamma_p;\mathcal H)$，
$J_W$ 取其弱导数。这里只需要证明
\begin{equation}\label{new:qualitative-spectrum}
 \lim_{r\downarrow0}\limsup_d\frac1p\mathbb E
 \operatorname{tr}\{r(pA_W+rI)^{-1}\}=0.
\end{equation}
以下直接使用全局匹配、$\mathbb E\|J_W\|_{\rm HS}^2\le C$ 和平稳 OU
余项 $\mathbb E\|R_t\|_{\mathcal H}^2\le Ct^{1+\delta}$；两支统一可取
$\delta=1/2$。

固定 $\mathcal H$ 的正交基 $(e_j)$，定义辅助 Gaussian 序列通道
\[
 Y_{\rho,j}=\sqrt{\rho p}\langle f(W),e_j\rangle+Z_j,
 \qquad K_d(\rho)=I(W;Y_\rho),\qquad Z_j\stackrel{\rm iid}{\sim}N(0,1).
\]
白噪声在序列空间定义，不视为 $\mathcal H$ 值变量。所有通道不等式先对有限
坐标成立，再取坐标数趋于无穷。

对 $W=\sqrt{1-t}V+\sqrt t\xi$，$V,\xi$ 独立标准 Gaussian，
给定 $V$ 后用均值 $f(V)$ 的 Gaussian 参考通道，得
\[
 \frac{K_d(\rho)}p\le\frac12\log\frac1t+
 \frac\rho2\mathbb E\|f(W)-f(V)\|^2.
\]
梯度矩和 OU 余项给最后的平方距离至多 $Ct$，故取 $t=\rho^{-1}$ 得
$K_d(\rho)/p\le\frac12\log\rho+C$，$\rho\ge2$。
令 $m_d(\rho)=\mathbb E\|f(W)-\mathbb E[f(W)\mid Y_\rho]\|^2$。
Gaussian I--MMSE 恒等式与 MMSE 单调性给
\[
 K_d(\rho)=\frac p2\int_0^\rho m_d(s)\,ds,
 \qquad m_d(\rho)\le\frac{2K_d(\rho)}{p\rho}
 \le C\frac{\log\rho}{\rho}.
\]
有限坐标时，取 $S=\sqrt p\,T_Mf(W)$，应用附录引理 \ref{info:immse}。
序列通道的恒等式由有限坐标互信息单调收敛以及支配收敛得到：
若 $T_M$ 为前 $M$ 坐标投影，则
$T_M\mathbb E[f(W)\mid Y_{1:M}]\to\mathbb E[f(W)\mid Y_{\rm all}]$
在 Hilbert $L^2$ 中收敛，且 MMSE 被 $\mathbb E\|f(W)\|^2$ 支配。

给定 $Y_\rho$ 独立抽取后验样本 $W,V$。平均后各自保持原 Gaussian
边缘，而函数距离平方的期望为 $2m_d(\rho)$。式 \eqref{eq:matching} 的全局匹配于是给
\[
 \mathfrak d_d(\rho):=\frac1p\mathbb E\min_{g\in\mathcal G_d}\|W-g(V)\|^2
 \le C\frac{\log\rho}{\rho}+e_d,\qquad e_d\to0.
\]
取可测最优标签 $G$。条件独立性、有限标签熵以及 Gaussian 最大熵给
\[
 h(W\mid Y_\rho)=h(W\mid Y_\rho,V)
 \le\log|\mathcal G_d|+\frac p2\log(2\pi e\mathfrak d_d(\rho)).
\]
此式亦可由以 $g(V)$ 为均值的有限 Gaussian 混合参考密度直接证明。
由于 $h(W)=\frac p2\log(2\pi e)$，对固定 $\rho\ge2$ 得
\begin{equation}\label{new:information-lower}
 \liminf_d\frac{K_d(\rho)}p
 \ge\frac12\log\rho-\frac12\log\log\rho-C.
\end{equation}

下面给出包含实际谱的辅助信息上界。条件于 $V$ 写
$f(W)-f(V)=\sqrt tJ_V\xi+R_t$。前 $M$ 个坐标中采用参考 Gaussian
协方差 $\Sigma_M=I_M+2\rho ptJ_MJ_M^*$，其中 $J_M=T_MJ_V$。
令 $L_M=\sqrt tJ_M\xi$、$R_M=T_MR_t$。由 $\|L_M+R_M\|_{\Sigma_M^{-1}}^2\le
2\|L_M\|_{\Sigma_M^{-1}}^2+2\|R_M\|^2$ 以及
$\mathbb E[L_ML_M^*\mid V]=tJ_MJ_M^*$，Gaussian KL 公式中的线性二次项
恰与 $\operatorname{tr}\Sigma_M^{-1}-M$ 抵消，得到
\[
 I(W;Y_{\rho,1:M}\mid V)
 \le\frac12\mathbb E\log\det\Sigma_M+\rho p\mathbb E\|R_t\|^2.
\]
令 $M\to\infty$ 并用 Sylvester 行列式公式，得
\begin{equation}\label{new:information-upper}
 \frac{K_d(\rho)}p\le\frac12\log\frac1t+
 \frac1{2p}\mathbb E\log\det(I_p+2\rho t\,pA_V)
 +C\rho t^{1+\delta}.
\end{equation}
这里不要求 $R_t$ 与 $\xi$ 独立。

记 $\nu_d$ 为 $pA_W$ 的期望经验谱测度，则
$\int\lambda\,d\nu_d=\mathbb E\operatorname{tr}A_W\le C$。
令 $q_d(r)=\nu_d([0,r])$，假设
$b=\lim_{r\downarrow0}\limsup_dq_d(r)>0$，固定 $0<\beta<b$。
取 $a\in(1/(1+\delta),1)$，并对每个固定大 $\rho$ 置
$t=\rho^{-a}$、$r_\rho=\rho^{-1}$。可选维数子序列使
$q=q_d(r_\rho)\ge\beta$。写 $Q=2\rho t$，Jensen 给
\[
 \int\log(1+Q\lambda)\,d\nu_d
 \le\log3+(1-q)\log\left(1+\frac{QC}{1-q}\right)
 \le(1-\beta)\log(1+QC)+C.
\]
沿此子序列，\eqref{new:information-upper} 的右侧至多
$\frac12[a+(1-\beta)(1-a)]\log\rho+C+o_\rho(1)$。
与 \eqref{new:information-lower} 比较，推出
$\beta(1-a)\log\rho\le\log\log\rho+C$，矛盾。因此 $b=0$。
最后利用
\[
 \int\frac r{\lambda+r}\,d\nu_d(\lambda)\le q_d(s)+r/s,
\]
依次取维数极限、$r\downarrow0$、$s\downarrow0$，即得
\eqref{new:qualitative-spectrum}。固定噪声方差 $\Delta$ 仅重标度 $r$。


因此实际 Fisher 的正则化亏损满足
\begin{equation}\label{eq:deficit}
 D_*(r):=\limsup_dD_d(r),\qquad
 \lim_{r\downarrow0}D_*(r)=0.
\end{equation}
不需要固定正谱边或谱尾的定量速率。

\subsection{预测 KL 下界}
附录 \ref{app:information} 对两支使用同一 OU 平滑测试场，证明
\begin{equation}\label{eq:lower-unified}
\epsilon_{\log,d,n}\ge
\frac{p[1-D_d(r)-C\sqrt{\dsm/r}]_+^2}
{2\{n(1+C\sqrt{\dsm/r})^2+Cp/(r^2\dsm)\}},
\qquad 0<r,\dsm\le1.
\end{equation}
取 $r=\alpha^{-1/8}$、$\dsm=\alpha^{-1/2}$。则
$\sqrt{\dsm/r}=\alpha^{-3/16}$、
$(\alpha r^2\dsm)^{-1}=\alpha^{-1/4}$；先令 $d\to\infty$，
再由 \eqref{eq:deficit} 令 $\alpha\to\infty$，得到
\begin{equation}\label{eq:lower-final}
\liminf_{\alpha\to\infty}\alpha\liminf_d\epsilon_{\log,d}(\alpha)\ge c/2.
\end{equation}

\section{粗信息界与后验函数收缩}
\newthought*{尖锐上界只需要一个非尖锐的初始收缩率。}
记 $I_{d,n}=I(\theta^0;\mathcal D_n)=I(W;\mathcal D_n)$。Bayes 定义和样本条件独立性给出
\begin{equation}\label{eq:increment}
\epsilon_{\log,d,n}=I_{d,n+1}-I_{d,n},\qquad
\epsilon_{\log,d,n+1}\le \epsilon_{\log,d,n}.
\end{equation}
第二式也可由条件熵递减看出：真实条件噪声熵不随数据量改变，预测熵随观察增加而减小。

\subsection{粗风险上界}
\newthought*{先暂时提供一个带噪参数观测。}
记 $\theta_t=\theta(V_t)$。令 $V_t=\sqrt{1-t}W+\sqrt tU$，$U$ 独立标准 Gaussian。Hermite OU 恒等式和梯度能量界给出
\[
\E\norm{F_\theta^{(L)}-F_{\theta_t}^{(L)}}_{L_X^2}^2\le Ct.
\]
链式法则与参考均值 $F_{\theta_t}^{(L)}$ 的 Gaussian 标签分布给出
\begin{equation}\begin{aligned}
I_{d,n}&\le I(W;V_t)+I(W;\mathcal D_n\mid V_t)\\
&\le\frac p2\log(1/t)+\frac n{2\Delta}
\E\norm{F_\theta^{(L)}-F_{\theta_t}^{(L)}}_{L_X^2}^2\\
&\le\frac p2\log(1/t)+Cnt.
\end{aligned}\end{equation}
取 $t=(1+2Cn/p)^{-1}$，得到
\[
I_{d,n}\le\frac p2\log(1+2Cn/p)+p/2.
\]
由 \eqref{eq:increment}，$n\epsilon_{\log,d,n}\le I_{d,n}$，所以
\begin{equation}\label{eq:coarse}
\limsup_d \epsilon_{\log,d}(u)\le C\frac{\log u}{u}
\quad(u\text{ 足够大}).
\end{equation}
此界只用先验梯度能量，未用待证首项。

\subsection{从 KL 到函数距离，再到匹配距离}
\newthought*{平均后验样本仍有精确的先验边缘。}
给定 $\mathcal D_n$ 独立抽取 $W,V$；平均数据后各自服从 $\Pi_d$。若 $B_X=F_\theta^{(L)}(X)-F_{\theta'}^{(L)}(X)$，Gaussian Hellinger 比较给出
\begin{equation}\label{eq:hellinger}
\E[1-e^{-B_X^2/(8\Delta)}]\le2\epsilon_{\log,d,n}.
\end{equation}
理由是 Hellinger 距离在 $L^2$ 平方根密度空间满足三角不等式；以后验预测密度为中间点，再用 $(x+y)^2\le2x^2+2y^2$，两个后验样本的平均平方距离至多四倍单样本平方距离，而后者由 KL 控制。Gaussian 两均值的亲和度恰为 $e^{-B_X^2/(8\Delta)}$。

对固定 $q>2$，先在 $B_X^2\le8\Delta$ 上使用
$1-e^{-x}\ge(1-e^{-1})x$，得到截断二阶矩 $\le C\epsilon_{\log,d,n}$。在补集上，\eqref{eq:hellinger} 控制概率，\eqref{eq:moments} 控制 $q$ 阶矩，Hölder 给出
\[
\E[B_X^2\boldsymbol1_{B_X^2>8\Delta}]
\le(\E|B_X|^q)^{2/q}\Pr(B_X^2>8\Delta)^{1-2/q}
\le C_q\epsilon_{\log,d,n}^{1-2/q}.
\]
因此
\[
\E\norm{F_\theta^{(L)}-F_{\theta'}^{(L)}}_{L_X^2}^2
\le C_q(\epsilon_{\log,d,n}+\epsilon_{\log,d,n}^{1-2/q}).
\]
将 \eqref{eq:matching} 应用于这一后验耦合，定义
\[
r_{d,n}=\frac1p\E\min_{g\in\mathcal G_d}\norm{W-gV}^2,
\qquad r_*(u)=\limsup_d r_{d,\lfloor ud^2\rfloor}.
\]
取 $q=2/\eta$，$0<\eta<1$，由 \eqref{eq:coarse} 得
\begin{equation}\label{eq:rate}
r_*(u)\le C_\eta\left(\frac{\log u}{u}\right)^{1-\eta}.
\end{equation}
这里只固定一次矩阶数，不令其随维度增长。

\section{一致 Gaussian 参考分布的样本块上界}
继续取 $V_t=\sqrt{1-t}W+\sqrt tU$，$0<t\le1/2$，反向写为
$W=\sqrt{1-t}V_t+\sqrt t\xi$，其中 $\xi$ 与 $V_t$ 独立标准 Gaussian。
记 $f_W=F_{\theta(W)}^{(L)}$、$a(x)=\nabla f_{V_t}(x)$、
$A=\E_Xa(X)a(X)^T$。已有 \eqref{eq:ou-rem}、\eqref{eq:moments} 给出
\[
 f_W(x)-f_{V_t}(x)=\sqrt t\,a(x)^T\xi+R_t(x),\qquad
 \E\|R_t\|_{L_X^2}^2\le Ct^{1+\delta},\quad
 \E\|a(X)\|^4\le C.
\]
两支统一可取 $\delta=1/2$。以下直接控制样本块的预测 KL，
不需要构造逐点均方误差估计器。

\subsection{相对熵链式法则消去参考模型误差}
给定 $V_t,X_{1:j}$，令 $\Phi_j$ 的第 $i$ 行为 $a(X_i)^T$，定义
\[
 Q_j=N\!\left((f_{V_t}(X_i))_{i\le j},S_j\right),
 \qquad S_j=\Delta I_j+t\Phi_j\Phi_j^T.
\]
这些参考分布由同一个虚拟系数 $\zeta\sim N(0,I_p)$ 的线性模型
$f_{V_t}(X_i)+\sqrt t\,a(X_i)^T\zeta+\sqrt\Delta Z_i$ 产生，
因此其样本边缘相容。输入采用真实的独立分布，虚拟系数不作为额外观察。
令 $P_j$ 为给定 $V_t$ 的真实数据分布，定义
\[
 C_j=\E\log\frac{p(Y_{1:j}\mid W,X_{1:j})}
 {q_j(Y_{1:j}\mid V_t,X_{1:j})},\qquad
 B_j=\E_{V_t}\KL(P_j\|Q_j).
\]
这里 $P_j,Q_j$ 均包含输入的相同边缘。相对熵分解和边缘相容性分别给
\[
 I(W;\mathcal D_j\mid V_t)=C_j-B_j,\qquad B_n-B_j\ge0\quad(n\ge j).
\]
其中第二式是相对熵链式法则，真实与参考的输入分布相同。
记知道 $V_t$ 的真实 Bayes 预测 KL 为 $\epsilon^V_{\log,d,k}$，因此
\begin{equation}\label{alt:block-entropy}
 \sum_{k=j}^{n-1}\epsilon^V_{\log,d,k}\le C_n-C_j.
\end{equation}

\subsection{交叉熵的非线性误差}
令 $L_j=\sqrt t\Phi_j\xi$、$R_j=(R_t(X_i))_{i\le j}$。
Gaussian 密度公式给出
\[
 C_j=\frac12\E\log\det(I_j+t\Phi_j\Phi_j^T/\Delta)
     +\E L_j^TS_j^{-1}R_j+\frac12\E R_j^TS_j^{-1}R_j.
\]
这里用到了精确抵消
\[
 \E[L_j^TS_j^{-1}L_j\mid V_t,X_{1:j}]
   +\Delta\tr S_j^{-1}=j.
\]
由于 $S_j\succeq\Delta I_j$，且线性模型的有效维数至多 $p$，
\[
 \E R_j^TS_j^{-1}R_j\le Cjt^{1+\delta},\qquad
 \E L_j^TS_j^{-1}L_j
 =\E\tr(t\Phi_j\Phi_j^TS_j^{-1})\le p.
\]
加权 Cauchy--Schwarz 与 Sylvester 恒等式遂给
\begin{equation}\label{alt:cross-entropy}
 \left|C_j-\frac12H_j\right|
 \le C\sqrt{pj t^{1+\delta}}+Cjt^{1+\delta},\qquad
 H_j=\E\log\det(I_p+t\Phi_j^T\Phi_j/\Delta).
\end{equation}
不要求 $R_j$ 与 $\xi$ 独立。

\subsection{四阶矩控制 log determinant 的 Jensen 亏损}
置 $\overline H_j=\E\log\det(I_p+jtA/\Delta)$。矩阵 $M\succeq I_p$ 上，
$-\log\det M$ 的 Hessian 满足
\[
 D^2(-\log\det)(M)[H,H]
 =\tr(M^{-1}HM^{-1}H)\le\|H\|_{\rm HS}^2.
\]
条件于 $V_t$，在矩阵
$I_p+(t/\Delta)\sum_{i\le j}a(X_i)a(X_i)^T$
与其均值间作二阶 Taylor 展开，再用条件独立性，得到
\[
 0\le\overline H_j-H_j
 \le\frac{jt^2}{2\Delta^2}\E\|a(X)a(X)^T-A\|_{\rm HS}^2
 \le Cjt^2.
\]
对 $n>j\ge1$，逐特征值使用
$(1+nt\lambda/\Delta)/(1+jt\lambda/\Delta)\le n/j$，即有
\[
 H_n-H_j\le\overline H_n-\overline H_j+Cjt^2
 \le p\log(n/j)+Cjt^2.
\]
结合 \eqref{alt:block-entropy}--\eqref{alt:cross-entropy}，得到有限维样本块界
\begin{equation}\label{alt:oracle-block}
 \sum_{k=j}^{n-1}\epsilon^V_{\log,d,k}
 \le\frac p2\log\frac nj+Cjt^2
      +C\sqrt{pn t^{1+\delta}}+Cnt^{1+\delta}.
\end{equation}
这一估计不需要实际或经验 Fisher 的谱下界。

\section{支付辅助信息并完成极限}
定义 $J_j(t)=I(W;V_t\mid\mathcal D_j)$。给定数据独立抽后验样本 $W'$，
用有限匹配族构造参考密度
\[
 Q_{W'}^{\rm par}=\frac1{|\mathcal G_d|}\sum_{g\in\mathcal G_d}
 N(\sqrt{1-t}\,gW',tI_p).
\]
任一混合分量都被 $|\mathcal G_d|Q_{W'}^{\rm par}$ 支配，Gaussian KL 公式给
\[
 \KL\!\left(N(\sqrt{1-t}W,tI_p)\middle\|Q_{W'}^{\rm par}\right)
 \le\log|\mathcal G_d|+\frac{1-t}{2t}\min_g\|W-gW'\|^2.
\]
相对熵分解及 $W'$ 的后验条件独立性因此给
\begin{equation}\label{alt:oracle-cost}
 J_j(t)\le\log|\mathcal G_d|+\frac{p r_{d,j}}{2t}.
\end{equation}
这里最优标签只用于上界，未提供给学习者。

由 Markov 链 $V_t-W-\mathcal D_n$ 和互信息链式法则，
\[
 \sum_{k=j}^{n-1}\epsilon_{\log,d,k}
 =\sum_{k=j}^{n-1}\epsilon^V_{\log,d,k}+J_j(t)-J_n(t).
\]
风险递减、$J_n(t)\ge0$ 及 \eqref{alt:oracle-block}--\eqref{alt:oracle-cost} 给
\begin{equation}\label{alt:full-block}\begin{aligned}
 (n-j)\epsilon_{\log,d,n}\le{}&
 \log|\mathcal G_d|+\frac{p r_{d,j}}{2t}+\frac p2\log\frac nj\\
 &+Cjt^2+C\sqrt{pn t^{1+\delta}}+Cnt^{1+\delta}.
\end{aligned}\end{equation}

固定 $\vartheta\in(0,1)$，取 $n=\lfloor\alpha d^2\rfloor$、
$j=\lfloor(1-\vartheta)\alpha d^2\rfloor$。所有小参数在维度极限中固定。
因为 $p/d^2\to c$、$\log|\mathcal G_d|=o(d^2)$，先令 $d\to\infty$ 得
\begin{equation}\label{alt:limit-upper}\begin{aligned}
 \alpha\limsup_d\epsilon_{\log,d}(\alpha)
 \le{}&\frac c{2\vartheta}\left\{
 \frac{r_*((1-\vartheta)\alpha)}t+\log\frac1{1-\vartheta}\right\}\\
 &+\frac C\vartheta\left\{\alpha t^2+
       \sqrt{\alpha t^{1+\delta}}+\alpha t^{1+\delta}\right\}.
\end{aligned}\end{equation}
在 \eqref{eq:rate} 中取 $\eta=1/8$，并取 $t=\alpha^{-3/4}$。
两支统一使用 $\delta=1/2$，则四个误差分别为
\[
 O_\vartheta(\alpha^{-1/8}(\log\alpha)^{7/8}),\qquad
 \alpha^{-1/2},\qquad\alpha^{-1/16},\qquad\alpha^{-1/8},
\]
全部趋零。先 $\alpha\to\infty$，再 $\vartheta\downarrow0$，由
$\vartheta^{-1}\log(1/(1-\vartheta))\to1$ 得
\[
 \limsup_{\alpha\to\infty}\alpha\limsup_d\epsilon_{\log,d}(\alpha)\le c/2.
\]
结合 \eqref{eq:lower-final}，并使用
\[
 \limsup_d|\alpha\epsilon_{\log,d}(\alpha)-c/2|
 =\max\{\alpha\limsup_d\epsilon_{\log,d}(\alpha)-c/2,
          c/2-\alpha\liminf_d\epsilon_{\log,d}(\alpha)\},
\]
即证明原定理。没有交换维度与样本比极限，也没有对累计信息渐近式求导。


\appendix
\section{固定深度网络的解析估计及分部积分的定义域}\label{an:appendix}\label{an:section}\label{app:local}
\newthought*{固定深度与比例宽度下，各阶矩和局部导数具有维数一致的控制。}
所有期望均针对独立标准高斯输入和权重；本节不把后验双副本当成高斯耦合。固定深度 $L$，令 $k_0=d$，且 $cd\le k_\ell\le Cd$，其中 $0<c<C<\infty$ 固定。采用
\[
 u^{(\ell)}=W^{(\ell)} h^{(\ell-1)}/\sqrt{k_{\ell-1}},\quad
 h^{(\ell)}=\sigma_\ell(u^{(\ell)}),\quad h^{(0)}=X,\qquad
 F=k_L^{-1/2}\boldsymbol1^Th^{(L)}.
\]
每层激活固定、全局 Lipschitz，且 $\mathbb E\sigma_\ell(G)=0$、$\mathbb E\sigma_\ell(G)^2=1$。常数可以依赖 $L,c,C$、激活以及指定的有限阶矩，不能依赖 $d$。后文的平滑常数还独立于 $0<\dsm\le1$；维数阈值也独立于该区间内的 $\dsm$。任何矩阶均先固定，再选择维数阈值。

\subsection{输出和隐藏范数的任意固定阶矩}
\begin{lemma}\label{an:moments}
令 $q_\ell=\|h^{(\ell)}\|^2/k_\ell$。每个固定 $r>0$ 均有
\[
 \mathbb E q_\ell^r\le C_r,\qquad
 \mathbb E|q_\ell-1|^r\le C_r d^{-r/2},\qquad
 \mathbb E|F|^r\le C_r.
\]
这些结论适用于混合使用上述激活的固定层数网络。此外，每个固定 $r>0$ 有 $\sup_{\ell,i,d}\mathbb E|h_i^{(\ell)}|^r<\infty$；任意固定个特征的固定次幂乘积具有一致有界的绝对期望，允许这些特征来自不同层或不同输入副本，只须各副本具有规定的高斯边缘分布。
\end{lemma}
\begin{proof}
写 $b(q)=\mathbb E\sigma(\sqrt qG)$ 和 $v(q)=\mathbb E\sigma(\sqrt qG)^2$。由 $|\sigma(z)|\le M|z|+M_0$，对全部 $q\ge0$ 有
\[
 \mathbb E|\sigma(\sqrt qG)|^s\le C_s(1+q^{s/2}),\quad
 |b(q)|\le C|q-1|,\quad |v(q)-1|\le C|q-1|. \tag{A1}\label{an:scalar}
\]
最后一式来自差平方公式及
$|\sqrt q-1|(\sqrt q+1)=|q-1|$；前一式用 $b(1)=0$。这些是全局估计，不剔除坏方差事件。

先记录样本均值的直接矩展开。条件独立、条件中心化的 $Y_i$ 若满足
$\mathbb E[|Y_i|^s\mid\mathcal F]\le C_s(1+q^{as})$，$2\le s\le2m$，则
\[
 \mathbb E\left[\left|N^{-1}\sum_iY_i\right|^{2m}\middle|\mathcal F\right]
 \le C_mN^{-m}(1+q^{2ma}). \tag{A2}\label{an:sample}
\]
事实上，展开 $2m$ 次幂；含只出现一次的指标的项为零。余项至多有 $m$ 个不同指标，有 $j$ 个不同指标的项数至多 $C_mN^j$；其条件矩绝对值至多 $C_m(1+q^{2ma})$。除以 $N^{2m}$ 即得所述界。

给定前层，当前层的各预激活恰为独立 $N(0,q_{\ell-1})$。Jensen 不等式先给出
$\mathbb E[q_\ell^{2m}\mid\mathcal F_{\ell-1}]\le C_m(1+q_{\ell-1}^{2m})$，故固定深度归纳得到正矩界。对
$Y_i=(h_i^{(\ell)})^2-v_\ell(q_{\ell-1})$
使用 \eqref{an:sample}，再使用 \eqref{an:scalar}，得到
\[
 \mathbb E|q_\ell-1|^{2m}\le C_md^{-m}
       +C_m\mathbb E|q_{\ell-1}-1|^{2m}.
\]
初值由 $q_0=d^{-1}\sum_iX_i^2$ 和同一矩展开给出 $O(d^{-m})$，故所有隐藏层均成立。
最后，条件于前 $L-1$ 层，精确分解为
\[
 F=\sqrt{k_L}\,b_L(q_{L-1})+
 k_L^{-1/2}\sum_i\{h_i^{(L)}-b_L(q_{L-1})\}.
\]
第二项的 $2m$ 阶条件矩由 \eqref{an:sample} 控制为 $C_m(1+q_{L-1}^m)$；第一项的 $2m$ 阶矩不超过
$C_mk_L^m\mathbb E|q_{L-1}-1|^{2m}=O(1)$。
取 $2m\ge r$ 并用矩的单调性完成输出矩证明。单个特征的矩由 \eqref{an:scalar} 对前层条件化得到；固定乘积的矩由 H\"older 不等式得到，不需各因子独立。特别地，全一读出的条件均值没有被丢弃。
\end{proof}

\subsection{前向、反向以及光滑 Hessian 的维数无关界}
令 $A_\ell=W^{(\ell)}/\sqrt{k_{\ell-1}}$，$B_\ell=\|A_\ell\|_{\rm op}$。以后 $P$ 表示 $1+\|X\|/\sqrt d$ 与有限个 $B_\ell$ 的非负固定次数多项式；存在第二组独立高斯矩阵时也允许其归一化算子范数。所有这样的 $P$ 的每个固定阶矩均一致有界。为核实这一点，在 $\mathbb R^m,\mathbb R^r$ 单位球面各取至多 $9^m,9^r$ 个点的 $1/4$ 网。对高斯 $m\times r$ 矩阵 $G$，双线性形式的网逼近给出
\[
 \|G\|_{\rm op}\le2\max_{u,v\text{ 属于网}}|u^TGv|,
 \quad
 \mathbb P\{\|G\|_{\rm op}^2\ge8[\log2+(m+r)\log9+s]\}\le e^{-s}.
\]
积分此尾界并使用有界宽度比，即得所需的全部矩。

\begin{lemma}\label{an:derivatives}
对 Lipschitz 激活，弱梯度满足 $\mathbb E\|\nabla_WF\|^r\le C_r$，每个固定 $r>0$ 均成立。若各激活还属于 $C^2$ 且二阶导数有界，则
\[
 \mathbb E\|\nabla_W^2F\|_{\rm HS}^2\le C.
\]
特别地，归一化 tanh 和中心化归一化 sigmoid 满足这些结论。
\end{lemma}
\begin{proof}
前向递推给出 $\|h^{(\ell)}\|/\sqrt{k_\ell}\le P$。定义
\[
 \beta_L=k_L^{-1/2}\boldsymbol1,\quad
 \lambda_\ell=\operatorname{diag}(\sigma_\ell'(u^{(\ell)}))\beta_\ell,\quad
 \beta_{\ell-1}=A_\ell^T\lambda_\ell.
\]
由有界斜率，$\|\beta_\ell\|+\|\lambda_\ell\|\le P$，且
\[
 \nabla_{W^{(\ell)}}F=\lambda_\ell(h^{(\ell-1)})^T/\sqrt{k_{\ell-1}},
 \qquad \|\nabla_WF\|^2\le P. \tag{A3}\label{an:gradient}
\]
网络在每个有限维空间局部 Lipschitz；链式公式在可微点成立，并代表弱导数。亦可逐层平滑激活，用上述界和弱紧致性识别同一弱梯度。

写 $J_{u,\ell},J_{h,\ell}$ 为对全部权重的 Jacobian。由
\[
 \|J_{u,\ell}\|_{\rm op}\le B_\ell\|J_{h,\ell-1}\|_{\rm op}
                  +\|h^{(\ell-1)}\|/\sqrt{k_{\ell-1}},\qquad
 \|J_{h,\ell}\|_{\rm op}\le M\|J_{u,\ell}\|_{\rm op},
\]
二者算子范数均不超过 $P$。计算图的二阶链式法则把输出 Hessian 精确写为各层的两类项之和。激活项为
\[
 J_{u,\ell}^T\operatorname{diag}(\beta_\ell\odot\sigma_\ell''(u^{(\ell)}))J_{u,\ell};
\]
它的 Hilbert--Schmidt 范数不超过
$\|J_{u,\ell}\|_{\rm op}^2\|\sigma_\ell''\|_\infty\|\beta_\ell\|$。
混合项的双线性形式为
\[
 k_{\ell-1}^{-1/2}\lambda_\ell^T
 \{U^{(\ell)}Dh^{(\ell-1)}[V]+V^{(\ell)}Dh^{(\ell-1)}[U]\}.
\]
其中一个方向的 Hilbert--Schmidt 范数平方恰为
$\|\lambda_\ell\|^2\|J_{h,\ell-1}\|_{\rm HS}^2/k_{\ell-1}$，
不超过 $\|\lambda_\ell\|^2\|J_{h,\ell-1}\|_{\rm op}^2$，因为该 Jacobian 的秩不超过 $k_{\ell-1}$。转置方向同理。对固定层数求和，再用 $P$ 的矩界即可。
\end{proof}
由两个引理及 Cauchy--Schwarz，对每个固定 $Q>0$ 还有
\[
 \mathbb E[(1+|F|^Q)\|\nabla F\|^2]\le C_Q,
 \qquad
 \mathbb E[\|\nabla F\|^2\boldsymbol1_{\{|F|>M\}}]\le C_QM^{-Q}. \tag{A4}\label{an:weightedmoments}
\]

\subsection{整个中心化分段线性支：负矩与加权门事件}
考虑固定实数 $a,b$ 的激活
\[
 \sigma(z)=az+b(|z|-c_0),\qquad c_0=\sqrt{2/\pi},
 \qquad a^2+b^2(1-2/\pi)=1. \tag{A5}\label{an:pl}
\]
当 $b=0$，这是已被激活类排除的线性函数 $\sigma(z)=\pm z$。以下取 $b\ne0$，但不要求 $a\ne0$，也不要求 $a+b$ 和 $a-b$ 均非零。因此偶函数支和一侧斜率为零的 ReLU 支均被包含。常数可以依赖固定的 $(a,b)$；不声称当 $b\to0$ 时一致。

\begin{lemma}\label{an:negative}
存在 $\eta_0,C_0>0$，使得对全部 $r\ge0$、$0<\eta\le\eta_0$，
\[
 \mathbb P\{\sigma(rG)^2\le\eta\}\le C_0\sqrt\eta.
\]
若 $Q=N^{-1}\sum_{i=1}^N\sigma(rG_i)^2$，则每个固定 $p>0$ 都存在 $N_0(p)$ 使
$\sup_{r\ge0,N\ge N_0(p)}\mathbb E Q^{-p}<\infty$。
\end{lemma}
\begin{proof}
记 $c=bc_0\ne0$。在两条半轴上，激活分别为 $s_\pm z-c$，其中 $s_\pm=a\pm b$。若斜率为零，该半轴取恒值 $-c$，当 $\sqrt\eta<|c|/2$ 时不会贡献小球事件。非零斜率的半轴上，小球事件属于以非零根 $c/s_\pm$ 为中心、长度 $2\sqrt\eta/|s_\pm|$ 的区间；无合法半轴根时直接舍去。缩小 $\eta_0$ 后，每个区间与零有固定正距离。对 $|z|\ge\delta>0$，密度 $r^{-1}\phi(z/r)$ 在全部 $r>0$ 上一致有界，因为令 $v=|z|/r$ 后它等于 $v\phi(v)/|z|$。$r=0$ 时事件为空。因此得小球界。

事件 $Q\le\eta$ 要求至少 $N/2$ 项不超过 $2\eta$；条件独立和子集计数给出
$\mathbb P(Q\le\eta)\le(C_1\eta^{1/4})^N$。
选更小的 $\eta_0$ 使 $C_1\eta_0^{1/4}<1/2$，则
\[
 \mathbb E Q^{-p}\le\eta_0^{-p}
       +pC_1^N\int_0^{\eta_0}s^{N/4-p-1}\,ds
 \le\eta_0^{-p}+\frac{p\eta_0^{-p}2^{-N}}{N/4-p}
\]
对 $N>4p$ 一致有界。逐层条件化立即证明全部隐藏范数及删去一个坐标后的平均平方范数具有固定阶负矩。输入范数的负矩则由
\[
 \mathbb E q_0^{-p}=\prod_{j=1}^p\frac d{d-2j}\le2^p\quad(d\ge4p)
\]
给出，此处先取整数 $p$，其余阶用矩单调性。
\end{proof}

下面给出避免把反传量误当独立变量的关键论证。固定一层 $\ell<L$，把下游范数多项式记为 $P_>$。条件于到该层为止的全部变量（允许同时给定另一个耦合表示 $h'$），当前网络的下游加权协方差是
\[
 M(h)=\mathbb E[P_>\beta_\ell\beta_\ell^T\mid h^{(\ell)}=h].
\]
被积函数只使用当前网络和下游高斯范数，并不使用 $h'$。下游第一张矩阵右乘任意正交矩阵不改变其高斯分布及范数；如果还有耦合高斯噪声矩阵，同时旋转它。因此 $M$ 正交等变。固定 $h$ 的正交稳定子推出
\[
 M(h)=c_\perp(I-hh^T/\|h\|^2)+c_\parallel hh^T/\|h\|^2.
\]
反向范数界完全不依赖输入 $h$，故 $\operatorname{tr}M(h)\le C$，从而对 $k=k_\ell\ge2$，
\[
 M(h)_{ii}\le C\left(k^{-1}+\frac{h_i^2}{\|h\|^2}\right). \tag{A6}\label{an:covariance}
\]
最后一层直接用 $\beta_{L,i}^2=1/k_L$。对任意前向界中的 $P$，逐单项分拆成前层、当前层和下游范数因子，即可应用此式。

处理当前层依赖时，删去第 $i$ 行，记为 $W^{(-i)}$；由
\[
 \|W/\sqrt{k_{\ell-1}}\|_{\rm op}
 \le\|W^{(-i)}/\sqrt{k_{\ell-1}}\|_{\rm op}
       +\|W_i\|/\sqrt{k_{\ell-1}},\qquad
 \|h\|^{-2}\le\left(\sum_{j\ne i}h_j^2\right)^{-1},
\]
可把当前范数多项式上界成删行变量多项式与归一化行长多项式的积。条件于前层，删行变量独立于第 $i$ 行。前一个引理和 H\"older 不等式给出
\[
 \mathbb E\left[\frac{P_{\rm del}}{\sum_{j\ne i}h_j^2}
                  \middle|\text{前层}\right]\le C/k. \tag{A7}\label{an:rowdelete}
\]
这使用任意所需的固定阶负矩和算子范数矩，并非以均值替代随机分母。

给定前层，任意成对高斯行的两个预激活可写成
$u=\sqrt qG$、$v=\sqrt r(\cos\omega G+\sin\omega H)$。
符号不同的事件在 $(G,H)$ 平面是总角度 $2\omega$ 的两个扇形。极坐标积分、其余独立高斯坐标的矩，以及 $h_i^2\le C(1+qG^2)$ 给出，对归一化行长的任意固定多项式 $Q$，
\[
 \mathbb E[Q\boldsymbol1_{uv<0}\mid\text{前层}]\le C_Q\omega,
 \quad
 \mathbb E[Qh_i^2\boldsymbol1_{uv<0}\mid\text{前层}]\le C_Q(1+q)\omega. \tag{A8}\label{an:angular}
\]
将 \eqref{an:covariance}--\eqref{an:angular} 合并并对 $i$ 求和，便得到加权门事件界；角度外的因子只留下前层范数多项式。没有任何一步假设 $\beta_{\ell,i}$ 独立于本层门。

\subsection{分段线性网络的直接 OU 梯度估计}
\begin{lemma}\label{an:gradient-ou}
令 $W'=\sqrt{1-t}W+\sqrt tZ$，$Z$ 为独立标准 Gaussian 权重，$0<t\le1/2$。
分段线性支满足
\[
 \mathbb E\|\nabla F(W',X)-\nabla F(W,X)\|^2\le C\sqrt t.
\]
常数与维数阈值均不依赖 $t$。
\end{lemma}
\begin{proof}
以下 $P$ 表示 $1+\|X\|/\sqrt d$ 及各层归一化 $W,Z$ 算子范数的任意固定非负系数多项式；这些量的全部固定阶矩一致有界。记撇号为 $W'$ 网络的量。由
$A'_\ell-A_\ell=(\sqrt{1-t}-1)A_\ell+\sqrt t\,Z^{(\ell)}/\sqrt{k_{\ell-1}}$
及 Lipschitz 前向递推，逐点有
\[
 \|A'_\ell-A_\ell\|_{\rm op}\le\sqrt t P,\qquad
 \frac{\|h'^{(\ell)}-h^{(\ell)}\|}{\sqrt{k_\ell}}\le\sqrt t P.
\]
给定前层 pair，本层各行 pair 相互独立，预激活相关系数为
$\sqrt{1-t}\langle h,h'\rangle/(\|h\|\|h'\|)$。令对应夹角为 $\omega$，则
\[
 \omega\le C\left(\sqrt t+
 \frac{\|h-h'\|}{\sqrt{\|h\|\|h'\|}}\right),\qquad
 \mathbb E[P\omega]\le C_P\sqrt t.
\]
末式由前向差界、引理 \ref{an:negative} 的两边缘负矩及 H\"older 得到；第一层 $h=h'=X$ 时直接成立。

记 $D_\ell=\operatorname{diag}(\sigma'(u^{(\ell)}))$。
我们先证明保留权重的门估计
\begin{equation}\label{an:pairedgate}
 \mathbb E\left[P\|(D'_\ell-D_\ell)\beta_\ell\|^2\right]
 \le C_P\sqrt t.
\end{equation}
将 $P$ 逐单项分为前层、本层和下游的因子。先给定前层 pair 及本层全部行，积分下游。
下游 $W,Z$ 同时右乘正交矩阵保持联合分布及范数，故单个原网络反传量仍满足
\eqref{an:covariance}；其条件协方差只依赖原表示 $h^{(\ell)}$，无需双网络交叉协方差。
本层同时删去 $W,Z$ 的第 $i$ 行。算子范数各由删行算子范数与归一化行长之和控制，
分母由 $\sum_{j\ne i}(h_j^{(\ell)})^2$ 控制。给定前层，删行变量与第 $i$ 行 pair 独立，
\eqref{an:rowdelete} 连同所需更高阶负矩使每行的删行因子至多 $C/k_\ell$。
第 $i$ 行剩余因子是两条 Gaussian 行长的固定多项式，可上界为
$C(1+\|(W_i,Z_i)\|/\sqrt{k_{\ell-1}})^D$，其中 $D$ 固定。
这一上界对联合行空间的正交旋转不变，且其固定阶矩与维数无关。\eqref{an:angular} 的二维扇区积分仍适用：
将 $(W_i,Z_i)$ 的两个预激活方向张成的平面与其正交补分开，径向矩一致有界，
得到该行门事件及其乘 $(h_i^{(\ell)})^2$ 的积分分别至多 $C\omega$、$C(1+q)\omega$。
两斜率差平方是该门事件的固定倍数。对 $i$ 求和，条件界因而至多
$CP_{<}(1+q)\omega$；再平均前层即得 \eqref{an:pairedgate}。终层直接用 $\beta_{L,i}^2=1/k_L$。

反传差满足精确递推
\[
 \lambda'_\ell-\lambda_\ell
 =D'_\ell(\beta'_\ell-\beta_\ell)+(D'_\ell-D_\ell)\beta_\ell,\qquad
 \beta'_{\ell-1}-\beta_{\ell-1}
 =A_\ell^T(\lambda'_\ell-\lambda_\ell)+(A'_\ell-A_\ell)^T\lambda'_\ell.
\]
从相同读出出发反向归纳，对任意上述 $P$，平方有限项之和并用
\eqref{an:pairedgate}，得
$\mathbb E[P(\|\beta'_\ell-\beta_\ell\|^2+\|\lambda'_\ell-\lambda_\ell\|^2)]\le C_P\sqrt t$。
乘上新的算子范数只改变 $P$，矩阵增量项的平方为 $O(t)$，故无需对门概率再取平方根。
最后将梯度块差写成
$[(\lambda'_\ell-\lambda_\ell)(h^{(\ell-1)})^T+
\lambda'_\ell(h'^{(\ell-1)}-h^{(\ell-1)})^T]/\sqrt{k_{\ell-1}}$，
使用上述带权界与前向差界并对固定层数求和，即得结论。
\end{proof}

\subsection{先验 OU 耦合的切线余项}
\begin{lemma}\label{an:ou}
令 $W(t)=\sqrt{1-t}\,W+\sqrt tZ$，其中 $Z$ 为独立标准高斯全参数向量，$0<t\le1/2$。则
\[
 \mathbb E|F(W(t),X)-F(W,X)-\sqrt t\,DF(W,X)[Z]|^2
 \le\begin{cases}Ct^2,&\text{光滑支},\\Ct^{3/2},&\text{分段线性支}.
 \end{cases}
\]
即 $O(t^{1+\delta})$ 中可分别取 $\delta=1$ 与 $\delta=1/2$。常数和维数阈值均不依赖 $t$。
\end{lemma}
\begin{proof}
条件于 $X=x$，将 $F$ 展开成正交归一化 Hermite 混沌，记第 $N$ 阶能量为 $c_N(x)$，并写 $a_t=\sqrt{1-t}$。生成函数给出 Mehler 恒等式
$\mathbb E H_\nu(W)H_\mu(W(t))=\boldsymbol1_{\nu=\mu}a_t^{|\nu|}$，
而 $\partial_iH_\nu=\sqrt{\nu_i}H_{\nu-e_i}$。于是
\[
 \mathbb E\|\nabla F(W(t))-\nabla F(W)\|^2
   =2\sum_{N\ge1}N(1-a_t^{N-1})c_N,
\]
\[
 \mathbb E|R_t|^2
  =\sum_{N\ge1}\{2(1-a_t^N)+tN-2tNa_t^{N-1}\}c_N.
\]
第二式中的交叉项亦可直接对 $Z$ 高斯分部积分得到。对 $N\ge2$，由凸性
$a_t^N=(1-t)^{N/2}\ge1-Nt/2$；对 $N=1$，$2(1-a_t)-t=(1-a_t)^2\le t^2$。故
\[
 \mathbb E|R_t|^2\le t^2\mathbb E\|\nabla F\|^2
             +t\mathbb E\|\nabla F(W(t))-\nabla F(W)\|^2. \tag{A12}\label{an:hermite}
\]
这些式子先对有限 Hermite 和成立，再在 Gaussian $W^{1,2}$ 中截断；引理 \ref{an:derivatives} 保证其定义域。
光滑支有 $\sum_NN(N-1)c_N=\mathbb E\|\nabla^2F\|_{\rm HS}^2\le C$，故梯度差平方为 $O(t)$。
分段线性支由引理 \ref{an:gradient-ou} 直接给出梯度差平方 $O(\sqrt t)$。
代入 \eqref{an:hermite} 完成证明。
这里保留完整 OU 余项；不能把径向漂移另行分离后假定其维数一致有界。
\end{proof}

\subsection{以参数 OU 半群平滑测试场}
\begin{lemma}\label{an:smoothing}
令 $(\mathsf P_sF)(W,X)=\mathbb E_ZF(e^{-s}W+\sqrt{1-e^{-2s}}Z,X)$，
$F_{\dsm}=\mathsf P_{\dsm^2}F$，$0<\dsm\le1$。
两支均满足 \eqref{eq:smoothing}，常数与维数阈值不依赖 $\dsm$。
\end{lemma}
\begin{proof}
记 $s=\dsm^2$、$c_N=\mathbb E_Xc_N(X)$。由 Hermite 梯度恒等式，
\[
 B(s):=\sum_{N\ge1}N(1-e^{-s(N-1)})c_N\le C\sqrt s,
 \qquad 0<s\le1.
\]
光滑支用 $1-e^{-x}\le x$ 及 Hessian 界得到更强的 $B(s)\le Cs$；
分段线性支用引理 \ref{an:gradient-ou}，其中 $t=1-e^{-2s}$。
当 $t>1/2$ 时用梯度二阶矩增大常数即可。OU 半群逐阶乘 $e^{-sN}$，故
\[
 \E\|\nabla F_{\dsm}-\nabla F\|^2
 =\sum_{N\ge1}N(1-e^{-sN})^2c_N
 \le Cs^2+CB(s)\le C\dsm.
\]
这里使用 $1-e^{-sN}\le(1-e^{-s})+(1-e^{-s(N-1)})$。
再由 $xe^{-2x}\le C(1-e^{-x})$，
\[
 \E\|\nabla^2F_{\dsm}\|_{\rm HS}^2
 =\sum_{N\ge2}N(N-1)e^{-2sN}c_N
 \le Cs^{-1}B(s)\le C/\dsm.
\]
最后 $\nabla\mathsf P_sF=e^{-s}\mathsf P_s\nabla F$，Jensen 和先验不变性给
$\E\|\nabla F_{\dsm}\|^4\le\E\|\nabla F\|^4\le C$。
计算先对有限 Hermite 和成立，再在 Gaussian Sobolev 范数下取极限。
\end{proof}

\subsection{Gaussian 散度、弱导数和截断的合法性}
\begin{lemma}\label{an:sobolev}
设 $g:\mathbb R^p\to\mathbb R^p$ 属于 Gaussian $W^{1,2}$，也允许值域带独立的外部 $L^2$ 变量。定义
$\delta g=\operatorname{div}g-W\cdot g$。则其闭扩张满足
\[
 \mathbb E(\delta g)^2
 =\mathbb E\|g\|^2+\mathbb E\sum_{i,j}(\partial_i g_j)(\partial_jg_i)
 \le\mathbb E\{\|g\|^2+\|Dg\|_{\rm HS}^2\}. \tag{A13}\label{an:divergence}
\]
\end{lemma}
\begin{proof}
先取紧支撑光滑场。对 Gaussian 密度逐坐标分部积分两次，利用
$\partial_i(W_jg_j)=\boldsymbol1_{i=j}g_j+W_j\partial_i g_j$，交叉项合并后正好得到等式；导数和中每项用 $2uv\le u^2+v^2$ 给出不等式。一般情形先用径向截断 $\chi_R$，其中 $0\le\chi_R\le1$、球内等于一、$\|\nabla\chi_R\|_\infty\le C/R$，再在紧集上卷积。Gaussian 密度在紧集上光滑且有正上下界，因此普通卷积逼近同时给出 Gaussian $W^{1,2}$ 逼近。尾部由 $g,Dg$ 的平方可积性控制，截断导数项由 $C\|g\|/R$ 控制。故有紧支撑光滑场序列在该 Sobolev 范数收敛；不等式保证其散度在 $L^2$ 中 Cauchy，并识别为分布散度。等式也由该收敛保留。对外部变量先用简单函数逼近并积分，证明不变。
\end{proof}

上述散度不等式用于附录 \ref{app:information} 的正则化测试场。
该场只依赖参数和独立测试输入；训练及辅助 Gaussian 观测的得分项另行处理。

\section{全局匹配所需的对称张量框架}\label{tg:section}
固定整数 $m\ge5$、$q\ge1$。设单位行 $w_1,\ldots,w_k$ 满足
$\max_{i\ne j}|\langle w_i,w_j\rangle|\le\delta$。对多重指标 $r\in\mathbb N^k$、$|r|=q$，定义
\[
 P_r=\operatorname{Sym}\bigotimes_jw_j^{\otimes mr_j},\qquad
 Z_{r,i}=\operatorname{Sym}\bigotimes_jw_j^{\otimes(mr_j-\mathbf1_{j=i})},\quad r_i>0.
\]
$\operatorname{Sym}$ 是全部槽位排列的平均，范数取继承的 Hilbert 张量范数。
令 $\mathcal F$、$\mathcal F^-$ 分别以 $P_r$、$Z_{r,i}$ 为列。则
\begin{equation}\label{tg:frame}
 cI\preceq\mathcal F^*\mathcal F\preceq CI,\qquad
 cI\preceq(\mathcal F^-)^*\mathcal F^-\preceq CI,
\end{equation}
只要 $\delta+\sum_{h=1}^q k^h\delta^{mh-2}=o(1)$；$c,C>0$ 只依赖 $m,q$。

证明使用精确配对公式
\[
 \left\langle\operatorname{Sym}\bigotimes_{a=1}^Ru_a,
 \operatorname{Sym}\bigotimes_{a=1}^Rv_a\right\rangle
 =\frac1{R!}\sum_{\pi\in\mathfrak S_R}\prod_a\langle u_a,v_{\pi(a)}\rangle.
\]
若列的标签重数为 $(n_j)$，对角中的无错配项为 $\prod_jn_j!/R!$，
它有只依赖固定总次数的正下界；其余对角项为 $O(\delta)$。
纯列的重数为 $mr$，删除列的重数为 $mr-e_i$。两族内部不同列具有不同重数，
后者因模 $m$ 的余数唯一确定 $i$。同一有限支撑内的非对角项至少含一个错配，
故为 $O(\delta)$，此类列的数量有固定上界。
若另一列含 $h\ge1$ 个新标签，它至少有 $mh-1$ 个槽必须与不同标签配对，
故内积至多 $C\delta^{mh-2}$；这样的列至多 $Ck^h$ 个。
因此每个 Gram 矩阵与上述正对角矩阵之差的绝对行和至多
$C(\delta+\sum_hk^h\delta^{mh-2})=o(1)$。对称矩阵的 Schur 界即得 \eqref{tg:frame}。

行长在固定正区间时，各列只乘有固定上下界的标量，结论不变。
对按输入宽度归一化的 Gaussian 行，行长一致趋于一且
$\delta=O(\sqrt{\log d/d})$、$k\asymp d$，故扰动趋零；事件补集概率可取任意固定负幂。

\section{原树系数核、fresh-coordinate 标记与 literal residual 的局部估计}
\label{fp:core-section}

本节只证明可逐次调用的局部命题。所有空间取实 Euclidean 空间及其 Hilbert
张量积；$\|\cdot\|_{\mathrm{HS}}$ 表示张量 Hilbert 范数，
$\operatorname{Flat}_{S|S^c}$ 表示按指定槽位分组的矩阵化。
各层宽度在 $[c d,Cd]$ 内，深度、次数、颜色数和操作次数均固定。
记 $\delta_d=L_d/\sqrt d$、$\rho_d=L_d/d$，其中 $L_d$ 是可随有限次
运算增大的固定对数幂。所有概率事件只在各自 Gaussian 边缘分布上建立；
事件成立后证明是确定性的，因此允许两个网络任意耦合。

\subsection{原树核的不变量及精确线性压缩}
\label{fp:core-invariant-section}

\paragraph{原树核的定义。}
从原网络的归一化读出向量开始，按原计算树自上而下展开。
每个当前模式代表一个原节点的指标，模式 $p$ 的一次展开为
\begin{equation}
 J e_a=\bigotimes_{t=1}^{q}u_a^{(t)},\qquad q\geq1,
 \qquad \|U_t\|_{\mathrm{op}}\leq C,
 \quad U_t=(u_a^{(t)})_a .
 \label{fp:core-replacement}
\end{equation}
允许一个因子为 $e_a$，表示保留原父指标；允许停止某个模式并保持其为
冻结叶。一个父指标的多个物理因子仍共享一个指标，不自动产生多个独立标记。
自适应的跨网络映射只能放在冻结叶上，且每个这样的映射输出恰好一个物理槽。
输出多个槽的非线性特征必须先作为树展开显式写入核。

\paragraph{引理（完整提升有界）。}
任何有限次上述展开的完整提升算子均有界。具体地，令
$\Delta_q e_a=e_a^{\otimes q}$，则
$J=(\bigotimes_tU_t^*)\Delta_q$，而 $\Delta_q$ 为等距嵌入，故
$\|J\|\leq\prod_t\|U_t\|$。
同样，若 $F_j:e_a\mapsto f_a^{(j)}$ 有界，则共享指标的乘积特征
$e_a\mapsto\bigotimes_jf_a^{(j)}$ 有界。
也可从 Gram 矩阵看出：其 Gram 为 $\mathop{\circ}_jF_j^*F_j$，
且对半正定矩阵 $X\preceq bI,Y\succeq0$，
$X\circ Y\preceq b\operatorname{diag}(Y)$。

\paragraph{引理（核的确定性闭合）。}
\label{fp:core-closure}
设 $E$ 有 $r$ 个模式，满足
\begin{equation}
 \|E\|_{\mathrm{HS}}\leq C,\qquad
 \max_{p,a}\|E(e_a\text{ 于模式 }p,\cdot)\|_{\mathrm{HS}}\leq\delta_d,
 \qquad
 \|\operatorname{Flat}_{S|S^c}E\|_{\mathrm{op}}\leq\delta_d
 \quad(\varnothing\ne S\ne[r]).
 \label{fp:core-invariants}
\end{equation}
除 $r=q=1$ 外，一次展开保持这些结论，常数可增大固定倍数。

\emph{证明。}
Hilbert 范数由完整提升有界立即得到。考虑新核的一个非平凡分割。
若新产生的整个块在同一侧，该分割就是旧核的非平凡分割两侧复合有界算子。
若新块跨越两侧，在两侧暂时保留正交父指标 $a$。中央算子按 $a$ 分块对角，
每块为旧切片 $E(e_a,\cdot)$ 的某个矩阵化，算子范数至多其 Hilbert 范数
$\delta_d$。再施加两侧的行映射即得结论，包含旧核只有一个模式的情形。

旧模式的坐标切片直接由有界提升控制。若固定新块第 $t$ 个因子的坐标 $j$，
且 $q\geq2$，先在其他新因子中保留正交父标签，得到平方范数
\[
 \sum_a |(U_t)_{aj}|^2\|E(e_a,\cdot)\|_{\mathrm{HS}}^2
 \leq\delta_d^2\|U_te_j\|^2\leq C\delta_d^2.
\]
剩余映射有界。若 $q=1,r\geq2$，新切片为
$E(U_1e_j,\cdot)$，由旧核第一矩阵化范数控制。
当 $r=q=1$，有界矩阵可能把分散向量集中于一个坐标，故不能作此推断。
\hfill$\square$

\paragraph{Gaussian 一元例外与原树实例。}
令 $A_\ell$ 为独立、按输入宽度平方根归一化的 Gaussian 矩阵。
任意固定连续积 $P=A_b\cdots A_a$ 满足
\begin{equation}
 \|P\|_{\mathrm{op}}\leq C,\qquad
 \max_{ij}|P_{ij}|\leq C_M\sqrt{\log d/d}
 \label{fp:core-product-event}
\end{equation}
且可使失败概率为 $O(d^{-M})$。算子范数由各层相乘得到；条件于
$Q=A_{b-1}\cdots A_a$ 后，$(A_bQ)_{ij}$ 是方差
$\|Qe_j\|^2/n_{b-1}\leq C/d$ 的 Gaussian 变量，尾界与有限并合即得第二式。
读出向量乘连续 Gaussian 积的坐标最大值也由条件于最后尚未暴露的
Gaussian 因子得到相同估计。于是首次分叉以前的 $r=q=1$ 链可单独处理；
首次分叉以后模式数不减，前一引理处理所有后续步骤。
冻结或恒等步骤无须概率估计。有限原树核全部满足
\eqref{fp:core-invariants}。特别强调：一元链内部不能插入自适应跨网络矩阵。

\paragraph{外部映射及重建。}
在每个核模式外施加有界的单槽映射，保持所有非平凡矩阵化界，但不必保持
坐标切片界。一次剪枝后，应从原树重新构造核：在新完成块的原父节点处
停止，并将其 metric 置于冻结叶。不能用“旧核任意收缩仍保持切片界”代替重建。
所有不同深度的冻结叶可同时存在，只要每条根至叶路径内部没有自适应 metric。

\paragraph{引理（线性槽的精确压缩）。}
\label{fp:core-linear}
令底层行向量为 $b_j$，父层矩阵为 $U=(u_{aj})$，并设
$z_a=\sum_ju_{aj}b_j$。一个父节点有 $t$ 个线性子槽以及次数
$q_1,\ldots,q_h\geq2$ 的非线性子槽，则其未对称化张量严格等于
\begin{equation}
 z_a^{\otimes t}\otimes
 \bigotimes_{\ell=1}^h\left(\sum_j u_{aj}b_j^{\otimes q_\ell}\right).
 \label{fp:core-linear-identity}
\end{equation}
\emph{证明。} 给每个有序子槽分配独立求和指标，再逐个使用多线性性；
线性槽的求和正好是 $z_a$，非线性槽保留原式。固定次数的对称化为槽排列的
平均，可在求和后施加。纯一元链因此是归一化矩阵的普通乘法，无额外
$\sqrt{\text{宽度}}$ 因子。式中的所有非线性因子仍共享 $a$，不得替换为独立阵列。
坐标掩码一般不与此求和交换；必须先在完整物理提升中证明掩码删除的误差为
$o(1)$，再恢复原张量并用本等式。\hfill$\square$

\subsection{fresh-coordinate 与真正的双原始标记}
\label{fp:core-fresh-section}

\paragraph{概率工具。}
若确定性矩阵 $A_j:\mathbb R^u\to\mathbb R^v$ 满足
$\max\{\|\sum_jA_j^*A_j\|,\|\sum_jA_jA_j^*\|\}\leq s^2$，则独立
标准 Gaussian 系数满足
\[
 \mathbb P\!\left(\left\|\sum_jg_jA_j\right\|>t\right)
 \leq2(u+v)\exp[-t^2/(2s^2)].
\]
这是有限维矩阵 Gaussian 级数界。其自伴膨胀为
$H_j=\left(\begin{smallmatrix}0&A_j^*\\ A_j&0\end{smallmatrix}\right)$，
$\|\sum_jH_j^2\|\leq s^2$；矩阵 Gaussian Laplace 界
$\mathbb E\operatorname{tr}\exp(\lambda_{\rm c}\sum_jg_jH_j)
\leq\operatorname{tr}\exp(\lambda_{\rm c}^2\sum_jH_j^2/2)$
配合 $\lambda_{\rm c}=t/s^2$ 得到上尾，对负号重复即得双侧界。
以下只用其维数为固定 $d$ 次幂时的 $O(s\sqrt{\log d})$ 结论。

\paragraph{引理（fresh-color 双标记）。}
\label{fp:core-fresh}
按原层同时展开一个核的所有未冻结模式。在该层 Gaussian 矩阵首次引入的
新颜色中选择一个模式 $p$，另选任意不同模式 $q$。设总模式数 $N\geq3$。
在裸核 $E$ 上，以任意高多项式概率有
\begin{equation}
 \max_j\sup_{\|v\|\leq1}
 \|E(e_j\text{ 于 }p,v\text{ 于 }q,\cdot)\|_{\mathrm{HS}}
 \leq\rho_d.
 \label{fp:core-fresh-slice}
\end{equation}
该层尚未展开下层后代；高层冻结颜色不参与新颜色的旋转。

\emph{证明。}
条件于一个满足 \eqref{fp:core-invariants} 的核，置
$A_jv=E(e_j,v,\cdot)$。直接按矩阵元素求和可知
\[
 \left\|\sum_jA_j^*A_j\right\|
 =\|\operatorname{Flat}_{\{q\}|\{q\}^c}E\|^2\leq\delta_d^2,
 \qquad
 \left\|\sum_jA_jA_j^*\right\|
 =\|\operatorname{Flat}_{\{p,q\}|\{p,q\}^c}E\|^2\leq\delta_d^2.
\]
第二分割非平凡恰好使用 $N\geq3$。前述级数界、
$\|g\|\geq c\sqrt d$ 的 Gaussian 尾界，给均匀球面向量 $u=g/\|g\|$
的收缩范数至多 $C_M\delta_d\sqrt{\log d/d}$。
对独立 Haar 矩阵的各列作并合，不要求列独立。
同一 Gaussian 层矩阵右乘该正交矩阵不改变分布；它同时旋转所有新颜色模式。
未固定模式上的旋转是等距算子，若 $q$ 也为新颜色，$v$ 上的上确界仍不变。
故 Haar 后的估计转回原坐标即为所需式。所有核和模式只有固定有限种。
\hfill$\square$

\paragraph{推论（碰撞恢复）。}
\label{fp:core-collision}
若新颜色出现至少两次，则总模式数至少三时，其碰撞扇区满足
\begin{equation}
 \|E_{\mathrm{coll}}\|_{\mathrm{HS}}^2
 \leq\sum_{p<q}\sum_j\|E(e_j\text{ 于 }p,e_j\text{ 于 }q,\cdot)\|^2
 \leq C d\rho_d^2=o(1).
 \label{fp:core-collision-bound}
\end{equation}
必须在新颜色首次引入时删去此张量，再让误差经过完整的有界提升。
这样最终误差仍为 $o(1)$，没有事后对教师指标求和产生的维数损失。
父指标的确定性复制不属于此处的碰撞；它始终是一个共享标签。

\paragraph{双父收缩与二阶非对角纤维。}
\label{fp:core-two-parent}
在一个未加掩码的原始切口中，设不同父模式 $p,q$ 的所选线性槽与教师行
收缩后给出 $\eta_i,\zeta_i$，其 $\ell^2$ 范数有界。模式 $p$ 是当前新颜色，
且删去所选槽后仍有非空残余；所有完整残余家族均为有界算子。
若切口总模式数 $R\ge3$，则完整剩余源张量 $S_i$ 满足
\[
 \|S_i\|_{\rm HS}\le C\rho_d.
\]
事实上，若 $q$ 也保留残余，先保留两个正交标签，平方范数至多
\[
 C\sum_{a,b}|\eta_{ia}\zeta_{ib}|^2
       \|E(e_a\text{ 于 }p,e_b\text{ 于 }q,\cdot)\|^2
 \le C\rho_d^2.
\]
若 $q$ 消失，先作其完整向量收缩，平方范数至多
\[
 C\sum_a|\eta_{ia}|^2
       \|E(e_a\text{ 于 }p,\zeta_i\text{ 于 }q,\cdot)\|^2
 \le C\rho_d^2.
\]
两式直接使用引理~\ref{fp:core-fresh}，均包括相同数值标签，无需碰撞掩码。
再施加残余家族、其余源提升和本教师叶剩余槽的收缩，范数只增大固定倍数。
与完整教师余切片配对并用
$\sum_i\|C(e_i,\cdot)\|\le C\sqrt d\|C\|$，得到
$O(\sqrt d\rho_d)\|C\|=o(1)$。
这里的两个标记是不同原始父位置；同一父节点的线性兄弟不产生第二标记。

二阶切口只需另用非对角元估计。原始同色二阶核 $D$ 是固定次数 Gaussian
多项式，且 $\mathbb E\|D\|_{\rm HS}^2\le C$；置换不变性给
$\mathbb E|D_{ab}|^2\le C/d^2$（$a\ne b$）。固定次数 Gaussian
超可积性给 $\|D_{ab}\|_{L^p}\le(p-1)^{D_0/2}\|D_{ab}\|_{L^2}$。
取 $p=C_M\log d$ 并对 $O(d^2)$ 个元并合，得到
$\max_{a\ne b}|D_{ab}|\le L_d/d$，失败概率 $O(d^{-M})$。
所需二阶矩由完整提升界及原 Gaussian 矩阵算子范数的固定阶矩得到。
该结论不控制二阶对角，也不控制二阶坐标--向量切片。

\paragraph{不能推广的标记。}
对 $v_a=d^{-1/2}$ 及归一化 Gaussian 行 $u_a$，核
$E_{a,j,k}=v_au_{aj}u_{ak}$ 具有有界 Hilbert 范数与密集最大元，然而
\[
 \left\|E\left(e_a,\frac{u_a}{\|u_a\|},\cdot\right)\right\|
 =d^{-1/2}\|u_a\|^2.
\]
两个可见后代不构成两个原始标记。总模式数为二时，
\eqref{fp:core-fresh-slice} 的第二方差界也不存在；例如
$E_{ab}=d_aU_{ab}$ 只有 $\|E\|_{\mathrm{op}}\leq C\|d\|_\infty$
可直接使用，不能声称所有坐标--向量切片为 $O(d^{-1})$。

\paragraph{推论（原始高阶碰撞的期望控制）。}
\label{fp:core-raw-collision-expectation}
设 $C_R$ 是固定深度、固定次数、零常数项 Gaussian 多项式网络的原始
齐次系数，$R\geq3$，允许固定有限个 profile 的和。则每个固定 $p<\infty$
均有 $\sup_d\mathbb E\|C_R\|_{\mathrm{HS}}^p<\infty$，并且
$\mathbb E\|\operatorname{Coll}_R C_R\|_{\mathrm{HS}}^2=o(1)$。
为证矩界，把每个 profile 写成有限个原树提升及槽排列的和。其范数至多
固定常数乘有限个原 Gaussian 矩阵算子范数的固定幂；归一化 Gaussian
矩阵在固定宽比下有维数一致的任意固定阶矩。因此 H\"older 不等式给所述
界，不需层间乘积各因子独立。对原始同色核，首个输入 Gaussian 层赋予
正交不变性，fresh 引理及碰撞推论给任意 $M$ 的事件 $\Omega_M$，其补集
概率 $O(d^{-M})$，而事件内碰撞平方范数至多 $L_d/d$。故
\[
 \mathbb E\|\operatorname{Coll}_R C_R\|^2
 \leq L_d/d+
 (\mathbb E\|C_R\|^4)^{1/2}\mathbb P(\Omega_M^c)^{1/2}=o(1).
\]
对称化只含固定有限个排列，有限 profile 和由三角不等式处理。

\subsection{完整余张量的统一局部收缩}
\label{fp:core-residual-section}

\paragraph{原始残余的界。}\label{fp:core-residual-flat}
设 $F_a=\sum_bU_{ab}Q_b$，其中 $U$ 是原 Gaussian 层或固定连续积，
$Q_b$ 为正次数原树 profile；不同高度的冻结叶允许任意有界自适应单槽映射。
完整家族 $e_a\mapsto F_a$ 有界。若 $F_a$ 至少有两个物理槽，则其每个
非平凡物理矩阵化至多 $C\delta_d$，同时对所有 $a$ 成立：以
$U_{a,:}$ 为读出重建裸核，式~\eqref{fp:core-product-event} 给读出的范数与
最大坐标，核闭合引理给矩阵化界，最后的单槽映射保持这些界。
家族有界性来自完整提升与 $U^*$ 的复合。单槽字面因子同样是有界家族，
不需要非平凡矩阵化估计。同指标乘积家族由对角嵌入保持有界。
以上不对自适应叶映射作独立性或坐标切片保持性假设。

\begin{lemma}[完整余张量的局部收缩]\label{fp:core-local-contraction}
\label{fp:core-split-hit}\label{fp:core-hybrid-Rge2}\label{fp:core-hybrid-R1}
固定一个教师叶及其所命中的源列，所有未消费物理槽始终保留。
教师叶含 $m\ge2$ 个物理槽。教师行矩阵和源完整家族有界；其他源列按各自标签分别完整提升，且不额外依赖所选列标签。教师系数为 $C$，源系数为 $D$（一模时记为 $d$）。
设该叶接触源列中的一个字面因子 $\ell_a$，其交叉矩阵
$G_{ia}=\langle a_i,\ell_a\rangle$ 满足
$\max|G|\le\kappa_d$ 及 $\sup_i\|G_{i,:}\|_2\le C$。
该列的其他因子为上段的原始残余或单槽字面因子。
则有下列三个完整配对界：
\begin{enumerate}
\item
若本教师叶还命中另一源列，则界为
$C\kappa_d\|C\|\|D\|$。
\item
若本叶全部位于该列，重建源核心有 $R\ge2$ 个模式，且
$\|\operatorname{Flat}_{1|\mathrm{rest}}D\|\le\delta_d$，则界为
$C\sqrt d\,\delta_d(\delta_d+\kappa_d)\|C\|$。
\item
若本叶全部位于该列，核心只有一个模式 $d_a$，
$\|d\|_\infty\le\kappa_d$，且除指定字面因子外至少有两个完整因子，
则界为 $C\sqrt d\,\kappa_d(\delta_d+\kappa_d)\|C\|$。
\end{enumerate}
各界在 $\delta_d,\kappa_d\le L_d/\sqrt d$ 时均为 $o(1)$。
一模仅一个残余因子不在本引理第三项内。
\end{lemma}
\begin{proof}
第一项先保留源标签 $a$，从其他列取本教师叶命中的全部 $t\ge1$ 个槽。
映射 $x\mapsto(\langle a_i^{\otimes t},x\rangle)_i$ 有界，因为其伴随
为 $(A^*)^{\otimes t}\Delta_t$。保留所得正交标签 $i$ 后，乘以
$G_{ia}$ 的范数至多 $\kappa_d$。随后在每个 $i$ 块内提升本列剩余家族，
收缩本叶余槽，并分析其余教师叶。这是从完整源系数到完整教师系数的
有界算子复合，故没有逐 $i$ 求和或部分迹损失。

以下考虑本叶全部位于该列。固定 $i$，按本叶如何消费各残余因子，分为
全消费、部分消费、全在外。全消费因子产生 $\ell^2$ 有界向量 $\eta$；
部分消费因子产生逐列范数至多 $C\delta_d$ 的有界家族；全在外家族仍有界。
记本列收缩后算子为 $B_i$。若存在部分消费因子，则
$\|B_i\|_{\rm HS}^2\le C\delta_d^2\sum_a|G_{ia}|^2$。
否则至少一个因子全消费，故
$\|B_i\|_{\rm HS}^2\le C\kappa_d^2\sum_a|\eta_a|^2$。
于是 $\|B_i\|_{\rm HS}\le C(\delta_d+\kappa_d)$。
对第二项，其他源模式的完整提升记为 $T$，Hilbert--Schmidt 理想不等式给
\[
 \|(B_i\otimes T)D\|_{\rm HS}
 \le\|T\|\,\|B_i\|_{\rm HS}
       \|\operatorname{Flat}_{1|\mathrm{rest}}D\|_{\rm op}
 \le C\delta_d(\delta_d+\kappa_d).
\]

对第三项，若存在部分消费且至少两个因子留有外部槽，残余乘积家族
的 Gram 是各家族 Gram 的 Schur 积。其中一个对角元至多 $C\delta_d^2$，
另一个至多 $CI$；由 $X\circ Y\preceq\|X\|\operatorname{diag}Y$
可知乘积家族范数至多 $C\delta_d$。系数
$d_aG_{ia}$ 乘全消费标量的 $\ell^2$ 范数至多 $C\kappa_d$。
若只一个因子留有外部槽且被部分消费，至少两个残余因子保证另有一个
全消费向量 $\eta$，于是三角不等式给
$\|S_i\|\le C\kappa_d\delta_d\sum_a|G_{ia}\eta_a|
\le C\kappa_d\delta_d$。
若没有部分消费而有外部槽，至少一个全消费向量把系数的 $\ell^2$ 范数
界为 $C\kappa_d^2$；若全部消费，至少两个全消费向量之积为
$\ell^1$ 有界，同样给 $C\kappa_d^2$。
因此第三项中完整源余张量满足
$\|S_i\|\le C\kappa_d(\delta_d+\kappa_d)$。
最后只需一次
$\sum_i\|C(e_i,\cdot)\|\le C\sqrt d\|C\|$，即得第二、三项。
\end{proof}

二阶源矩阵若先限制到非对角，可分别处理 $D$ 与 $\operatorname{diag}D$，
因为后者的算子范数及 Hilbert 范数不增；这里没有将其对角宣称为小量。
其他掩码不能据此处理。一模单残余的永久标记论证仍单独保留。

\paragraph{永久根指标与确切边界。}
\label{fp:core-spectator}
若一模例外仅有一个 residual 因子，不能使用前引理。
可把原根指标 $a$ 永久保留，以 $E_{ab}=d_aU_{ab}$ 开始，只展开后代模式。
由核闭合引理，所有重建快照都有至少两个模式，全部非平凡矩阵化为
$O(\delta_d)$；根 literal 作为一个有界外部单槽映射。
因此后代 hybrid 若确有分离的完整特征表示，可用
引理~\ref{fp:core-hybrid-Rge2}。当两个 fresh 活跃模式存在时，永久根指标
保证总模式数至少三，亦可用引理~\ref{fp:core-fresh}。
这不是删除根指标或对其取迹的许可。

\paragraph{低阶情形。}
二模双父状态不能套用
$d^{-1}$ 标记界，一模仅一个 residual 因子不能套用一模乘积引理。
特别是仅知某个已消费行核 $T$ 很小并不够：取
$C_{ir}=\delta_{ir}/\sqrt d,T_{ib}=\delta_{ib}/\sqrt d,F_{br}=\delta_{br}$，
则 $\|C\|_{\mathrm{HS}}=1$、$\|C\|_{\mathrm{op}}=\|T\|_{\mathrm{op}}=d^{-1/2}$、
$\|F\|=1$，但 $\sum_{i,b,r}C_{ir}T_{ib}F_{br}=1$。
这说明丢失原系数之后的部分迹可恢复维数因子；上文完整源张量证明避免了
这一操作。完整历史的因子分解将在第~\ref{fp:history} 节中逐步构造。

\section{实际激活的固定输入混沌与原始正次多项式闭合}
\label{fp:actual-section}\label{app:closure}

本节只研究同一网络先验下的值误差。令 $k_0=d$，$L$ 固定，且
$c d\leq k_\ell\leq C d$。写 $A_\ell=W^{(\ell)}/\sqrt{k_{\ell-1}}$，
$W^{(\ell)}$ 的元素独立服从 $N(0,1)$，各层独立；$X\sim N(0,I_d)$ 独立于权重。
定义
\[
 h^{(0)}=X,\quad h^{(\ell)}=\sigma(A_\ell h^{(\ell-1)}),\quad
 q_\ell=k_\ell^{-1}\|h^{(\ell)}\|^2,\qquad
 F_L:=F_{\theta(W)}^{(L)}=k_L^{-1/2}\sum_i h_i^{(L)}.
\]
激活满足 $\mathbb E\sigma(G)=0$、$\mathbb E\sigma(G)^2=1$，并取以下之一：
有界光滑奇函数（每个固定阶导数有界），或
\begin{equation}\label{fp:actual-pl}
 \sigma(z)=a z+b(|z|-u),\qquad u=\sqrt{2/\pi},\qquad
 a^2+b^2(1-u^2)=1.
\end{equation}
记 $\psi_j=\mathrm{He}_j/\sqrt{j!}$，
$a_j=\mathbb E[\sigma(G)\psi_j(G)]$，
$b_j=\mathbb E[(\sigma(G)^2-1)\psi_j(G)]$。
$\Pi_N$ 表示输入 $X$ 的第 $N$ 阶正交混沌投影；对称张量约定为
\begin{equation}\label{fp:actual-isometry}
 \mathbb E_X I_N(T)I_N(S)=N!\langle T,S\rangle.
\end{equation}
特别地，下文的 ``形式次数为 $N$ 的系数再取 $I_N$''，不等于先计算整个
原始多项式函数再作 $\Pi_N$。后者会包含更高形式次数的收缩项。

\paragraph{局部矩估计。}
以下估计均只涉及单网络边缘分布。
这些结论分别由定理~\ref{an:moments}、条件行估计~\ref{an:scalar}
及~\ref{an:derivatives} 中的 Gaussian 矩阵网格尾界证明。
对固定层数、固定 $p<\infty$，有
\begin{equation}\label{fp:actual-local-interface}
 \|q_\ell-1\|_{L^p(W,X)}\leq C_p d^{-1/2},\quad
 \sup_{d,\ell,i}\mathbb E|h_i^{(\ell)}|^p<\infty,\quad
 \sup_d\mathbb E\|A_\ell\|_{\rm op}^p<\infty.
\end{equation}
同样的固定乘积矩界适用于任意固定多个、允许重复标签的实际特征。
从第一式和任意大的 $p$，$q_\ell\notin[1/2,2]$ 的概率可取任意多项式衰减。
首层单位方向 $v_i=A_{1,i:}/\|A_{1,i:}\|$、半径 $s_i=\|A_{1,i:}\|$
同时满足
\begin{equation}\label{fp:actual-row-event}
 \max_i|s_i-1|+\max_{i\ne j}|\langle v_i,v_j\rangle|
 \leq C_M\sqrt{\log d/d},\qquad \|A_1\|_{\rm op}\leq C_M
\end{equation}
的失败概率至多 $d^{-M}$。这些结论也只涉及单个网络先验。
下面将从这些局部界给出截断、集体均值和有限轮恢复的证明；原始系数的
高阶碰撞界将在使用处精确列明。

\subsection{有限个实际脊函数的完整输入投影}
\label{fp:actual-finite-support}
固定 $t,N$，设 $v_1,\ldots,v_t$ 为单位向量，$G=(\langle v_i,v_j\rangle)$，
$\|G-I\|\leq\eta$，$\max_i|s_i-1|\leq\delta$。则
\begin{equation}\label{fp:actual-ridge-projection}
 \left\|\Pi_N\prod_{i=1}^t\sigma(s_iv_i\cdot X)
 -\sum_{\substack{n_i\geq1\\\sum_i n_i=N}}
 \frac{\prod_i a_{n_i}}{\sqrt{\prod_i n_i!}}
 I_N\!\left(\operatorname{Sym}\bigotimes_i(s_iv_i)^{\otimes n_i}\right)
 \right\|_{L^2_X}\leq C_{t,N}(\eta+\delta).
\end{equation}
$t>N$ 时右侧的近似和为空。

证明如下。所有相关张量位于 $\operatorname{Sym}^N\operatorname{span}\{v_i\}$。
用 $n_i\geq0$、$\sum n_i=N$ 的
$B_n=I_N(\operatorname{Sym}\bigotimes_i v_i^{\otimes n_i})$ 作基。
其 Gram 矩阵是 $G$ 的多项式，且在 $G=I$ 附近一致可逆：从标准正交基
到 $v_i$ 的线性映射有一致有界的最大、最小奇异值，其固定对称张量幂也如此。
投影的系数向量为此 Gram 矩阵的逆乘
$\mathbb E[\prod_i\sigma(s_iZ_i) B_n(Z;G)]$，其中 $Z\sim N(0,G)$。
在 $G$ 的固定紧邻域内，对 Gaussian 密度关于协方差求导，得到一致可积的
多项式乘 Gaussian 密度；$B_n$ 的系数也多项式依赖于 $G$。
激活至多线性增长，故这些有限个积分关于 $G$ 局部 Lipschitz。
关于 $s_i$ 的 Lipschitz 性来自 $\sigma$ 的 Lipschitz 性和 Gaussian 加权矩。
在 $G=I,s_i=1$，独立性给出系数 $\prod_i a_{n_i}/\sqrt{\prod_i n_i!}$；
$a_0=0$ 去掉所有零次数。最后把 $v_i^{\otimes n_i}$ 改回
$(s_iv_i)^{\otimes n_i}$ 只产生 $O(\delta)$。有限维范数等价证明
\eqref{fp:actual-ridge-projection}。这个证明对完整非多项式乘积积分，
因此已经包括高叶 Hermite 次数和非正交源方向造成的一切低阶收缩。

\subsection{实际正根展开的有限化：截断顺序与误差}
\label{fp:actual-positive}
给定前层特征 $h\in\mathbb R^k$，$q=\|h\|^2/k$，当前行预激活为
$W_i\cdot h/\sqrt k$。对 $q>0$ 写
\[
 \sigma(\sqrt q G)=\sum_{r\geq0}c_r(q)\psi_r(G),\qquad
 c_r(q)=\mathbb E\sigma(\sqrt qG)\psi_r(G).
\]
多元 Hermite 加法公式为
\begin{equation}\label{fp:actual-multivariate}
 \psi_r(W_i\cdot h/\|h\|)=
 \sum_{|\alpha|=r}\sqrt{\frac{r!}{\alpha!}}
 (h/\|h\|)^\alpha\prod_j\psi_{\alpha_j}(W_{ij}).
\end{equation}
$q=0$ 按连续的实际随机变量取值，不引入 $q^{-1}$ 的矩。

对本节出现的有限函数集 $f\in\{\sigma,\chi=\sigma^2-1\}$，其中 $\chi$
只用于分段线性分支的顶根，均有 $|f(x)|\le C(1+|x|^2)$。
映射 $q\mapsto f(\sqrt q\,\cdot)$ 在 $[1/2,2]$ 上为 $L^2(\gamma)$ 连续，
故其 Hermite 投影尾在此紧区间上一致趋零。在区间外，投影的 $L^2$ 范数
不超过完整函数，条件二阶矩至多 $C(1+q^2)$；式~\eqref{fp:actual-local-interface}
给任意高阶的 $q$ 矩和区间外概率，从而区间外误差趋零。
固定多个不同父行的乘积也同样可截断：条件于前层，行独立，逐因子写差后，
每项的条件平方期望等于一个尾的平方期望乘其余因子的条件二阶矩。
在紧区间上用一致尾界，区间外用上述多项式矩界即可，不需要单行尾的高阶
$L^p$ 收敛。根至叶依次选择截断，每次父行数已经固定。

固定保留次数 $r$，令 $c_r^f(q)=\mathbb E f(\sqrt qG)\psi_r(G)$。
将积分变量换成 $y=\sqrt qG$ 后，可对 Gaussian 密度和
$\psi_r(y/\sqrt q)$ 求导；在 $[1/2,2]$ 上导数由可积 Gaussian 多项式控制。
故 $c_r^f(q)q^{-r/2}$ 局部 Lipschitz，可替换为 $c_r^f(1)$，误差由
$\|q-1\|_{L^p}=O(d^{-1/2})$ 支付。区间外，精确有限分量的条件固定阶矩
由 Gaussian Hermite 矩界控制，系数是 $c_r^f(q)$，没有逆方差；替换后的
固定多项式由实际特征矩控制。H\"older 不等式与任意高阶区间外概率使误差
趋零。此论证从不使用负阶 $q$ 矩，也不要求分段线性 $f$ 可微。

接着在每个界面删去重复子标签，\emph{此时系数仍为正规序 Hermite，子特征仍为实际函数}。
有固定 $T$ 个子槽时，碰撞赋值最多 $C_T k^{T-1}$ 个，平方归一化为
$k^{-T}$。对当前权重的多指标 Hermite 正交性以及重复实际特征的固定乘积矩，
给出整个碰撞误差平方期望 $O_T(k^{-1})$。
不同排序若描述同一边集合，仅有依赖于固定 $T$ 的重数，应先合并再用正交性。
只在删去碰撞后，每个权重变量出现一次，正规序系数才能改写成普通边变量乘积。
对分段线性分支，出现在正根固定乘积\emph{内部}的零次项
$c_0(q)=bu(\sqrt q-1)$ 有任意固定 $L^p$ 范数 $O(d^{-1/2})$，
故由同一乘积估计去掉；顶根零次项另行保留。

为说明这些局部误差不会被宽度放大，将某阶段的整个误差写为
\begin{equation}\label{fp:actual-phase}
 \Delta=\sum_g\beta_g\Psi_g(W_{>\ell}) E_g(W_{\leq\ell},X),
 \qquad \sum_g\beta_g^2\leq C_{\rm fixed}.
\end{equation}
这里 $\Psi_g$ 是已经处理的不同平方自由边集合的乘积；相同边集合的描述
先合并。每条边带一个输入宽度的 $-1/2$ 次幂，根读出带 $k_L^{-1/2}$。
对根标签及每个自由下行子标签求和，恰好分别消去这些平方归一化；
只留下固定图形、阶数、排序的组合常数，证明第二个不等式。
$\Psi_g$ 只依赖更高独立 Gaussian 层，故
\begin{equation}\label{fp:actual-phase-bound}
 \mathbb E\|\Delta\|_{L^2_X}^2
 \leq C_{\rm fixed}\sup_g\mathbb E\|E_g\|_{L^2_X}^2.
\end{equation}
如果某误差共享同一高层相位，先将其整个低层乘积作为 $E_g$，不把共享相位
误当独立。于是固定 $\ecut>0$ 后，依次选择有限的根至叶截断，
得到有限个全界面标签互异图值之和 $F_{+,\ecut}$，且
\begin{equation}\label{fp:actual-finite-error}
 \limsup_d\mathbb E\|F_+-F_{+,\ecut}\|_{L^2_X}^2
 \leq C_L\ecut^2.
\end{equation}
所有归一化、径向替换和碰撞误差均在单网络边缘分布下支付。

最后保留首层实际脊函数，应用\eqref{fp:actual-ridge-projection}。
对每个固定支持大小，\eqref{fp:actual-row-event} 使其误差趋零；利用
\eqref{fp:actual-phase-bound} 聚合而不是对所有支持作三角不等式。
在首层好事件补集上，固定次数多项式矩及实际固定乘积矩至多有固定多项式维数增长，
可选足够大的失败指数使贡献趋零。
若图中任何节点度数大于 $N$，它有超过 $N$ 个不同子树；每棵子树有正次后代，
因而有超过 $N$ 个互异首层叶。其主项由\eqref{fp:actual-ridge-projection}为空。
所以所有截断均至少 $N$ 后，非空主图集合不再改变。
固定截断先令 $d\to\infty$，再令 $\ecut\to0$，得到\emph{同一个固定有限族}，
而非随精度增大的最终表达式。

\subsection{顶根递推与固定形式系数}
\label{fp:actual-mean}\label{fp:actual-masked-formula}
固定 $N>2$，定义零常数多项式
\begin{equation}\label{fp:actual-PQ}
 P_N(z)=\sum_{j=1}^N\frac{a_j}{\sqrt{j!}}z^j,
 \qquad Q_N(z)=\sum_{j=1}^N\frac{b_j}{\sqrt{j!}}z^j.
\end{equation}
用原始矩阵及归一化读出组成深度 $t$ 的形式多项式网络，根多项式为 $A$，
各下层多项式为 $p_1,\ldots,p_{t-1}$。将形式 $N$ 次对称系数取 $I_N$，
记为 $R_{t,N}[A;p]$；在每一神经元界面对全部子槽施加标签全局互异掩码，所得量记为
$J_{t,N}[A;p]$。输入物理坐标不加掩码，故 $J_{1,N}=R_{1,N}$。
下文 $;P_N$ 表示所有下层均用 $P_N$。固定辅助标记取系数亦包含在此记号内。

令 $F_t^f=k_t^{-1/2}\sum_i f((A_th^{(t-1)})_i)$，并用 $+$ 标记关于根
权重的正 Hermite 阶部分。前述有限化、相位正交与固定 $N$ 稳定性给
\begin{equation}\label{fp:actual-positive-local}
 \Pi_NF_{t,+}^{f}=J_{t,N}[H_f;P_N]+o_{L^2(W,X)}(1),
 \qquad H_\sigma=P_N,\quad H_\chi=Q_N.
\end{equation}
这里 $\chi$ 只用于分段线性分支。深度一没有高层相位，须直接检查：
$f=\sigma,\chi$ 满足
$|f(sx)-f(tx)|\le C|s-t|(|x|+x^2)$（$s,t\in[1/2,2]$），故
$\mathbb E[f(sG)\psi_N(G)]$ 对 $s$ 局部 Lipschitz。以其 $s=1$ 值乘
$s^N$ 替换的逐行误差为 $O(\max_i|s_i-1|)$；映射
$e_i\mapsto v_i^{\otimes N}$ 有界，因为其 Gram 绝对行和至多
$\sum_j|\langle v_i,v_j\rangle|^2\le\|V\|_{\rm op}^2$。
归一化读出范数为一，故浅层脊和误差也趋零；坏事件由固定矩与任意高尾概率支付。

分段线性支有下列对所有 $q\ge0$ 成立的精确恒等式：
\begin{equation}\label{fp:actual-root-mean}
 \mathbb E f(\sqrt qG)=\lambda_f(q-1)+c_f(\sqrt q-1)^2,
 \quad
 (\lambda_\sigma,c_\sigma)=(bu/2,-bu/2),\qquad
 (\lambda_\chi,c_\chi)=(1,b^2u^2).
\end{equation}
由 $|\sqrt q-1|\le|q-1|$ 及四阶局部矩，乘 $\sqrt{k_t}$ 后的二次余项
为 $O_{L^2}(d^{-1/2})$。又 $F_{t-1}^\chi=\sqrt{k_{t-1}}(q_{t-1}-1)$，所以
\begin{equation}\label{fp:actual-root-recursion}
 \Pi_NF_t^f=J_{t,N}[H_f;P_N]
 +\lambda_f\sqrt{k_t/k_{t-1}}\,\Pi_NF_{t-1}^\chi+o_{L^2}(1),
 \qquad t\ge2.
\end{equation}
$t=1$ 时直接有 $\Pi_NF_1^f=R_{1,N}[H_f]+o_{L^2}(1)$：条件均值的线性项
为 $\lambda_f\sqrt{k_1}(q_0-1)$，它属于输入二阶混沌，故被 $N>2$ 投影
精确消去；二次余项仍为 $o(1)$。这也证明
\eqref{fp:actual-positive-local} 的深度一版本。奇分支所有条件均值为零，
只需 $f=\sigma$，递推中的第二项不存在。

\subsection{二阶对角与有限降深度恢复}
\label{fp:actual-quadratic}\label{fp:actual-telescope}
原始碰撞推论~\ref{fp:core-raw-collision-expectation} 给出：固定零常数
多项式网络的全原始 $r\ge3$ 次系数满足
\begin{equation}\label{fp:actual-raw-interface}
 \mathbb E\|\operatorname{Coll}_rD_r\|_{\rm HS}^2=o(1),\qquad
 \sup_d\mathbb E\|D_r\|_{\rm HS}^p<\infty\quad(p<\infty).
\end{equation}
此结论不应用于带掩码系数。二阶对角须单独计算。

设输入宽度为 $k$ 的深度 $t$ 上部网络各层多项式为
$p_j(z)=\sum_{r\ge1}g_{j,r}z^r$。其全二次系数满足
\begin{equation}\label{fp:actual-diag}
 \mathbb E\left\|\operatorname{Diag}D_2-
 \lambda I_k/\sqrt k\right\|_{\rm HS}^2=O(k^{-1}),\qquad
 \lambda=g_{t,2}\prod_{j<t}g_{j,1}^{2}\sqrt{k_t/k}.
\end{equation}
正次网络的二次项恰有一个二次节点。若在第 $j$ 层，去掉固定标量后的系数为
$D_2^{(j)}=B_j^T\operatorname{diag}(v_j)B_j$，其中
$B_j=A_j\cdots A_1$、$v_j=A_{j+1}^T\cdots A_t^T\mathbf1/\sqrt{k_t}$。
写 $B_j=C_jA_1$、$H_j=C_j^T\operatorname{diag}(v_j)C_j$。
条件于 $H_j$，$A_1$ 各列为独立 $N(0,I/k)$，故
\[
 \mathbb E(D_2^{(j)})_{ii}=\operatorname{tr}H_j/k,
 \qquad
 \mathbb E\sum_i|(D_2^{(j)})_{ii}-\operatorname{tr}H_j/k|^2
 =2\|H_j\|_{\rm HS}^2/k.
\]
$\|H_j\|_{\rm HS}\le\|C_j\|_{\rm op}^2\|v_j\|$ 与固定矩界控制后式。
$j<t$ 时 $v_j$ 独立于 $C_j$，均值零、协方差 $I/k_j$，从而
$\mathbb E|\operatorname{tr}H_j|^2=k_j^{-1}\mathbb E\sum_a\|C_{j,a:}\|^4=O(1)$，
标量对角平方范数亦为 $O(k^{-1})$。$j=t$ 时
$\operatorname{tr}H_t=\|C_t\|_F^2/\sqrt{k_t}$，均值为 $\sqrt{k_t}$，
且 $\operatorname{Var}\|C_t\|_F^2=O(1)$：沿固定长度乘积逐层使用
\[
 \mathbb E(\|AC\|_F^2\mid C)=\frac{k_{\rm new}}{k_{\rm old}}\|C\|_F^2,
 \qquad
 \operatorname{Var}(\|AC\|_F^2\mid C)
 =\frac{2k_{\rm new}}{k_{\rm old}^2}\|CC^T\|_{\rm HS}^2=O_{L^1}(1)
\]
及全方差公式即可。因此只有顶层二次节点给确定主对角，恢复固定标量
得到\eqref{fp:actual-diag}；辅助标记取系数不改变证明。

\begin{lemma}[有限降深度恢复]\label{fp:actual-restoration}
对任意固定零常数根多项式 $A$，有
\begin{equation}\label{fp:actual-triangular}
 J_{t,N}[A;p]=R_{t,N}[A;p]
 -\sum_{s=1}^{t-1}\lambda_{t,s}(A)\,
 J_{s,N}[p_s^2;p_{<s}]+o_{L^2}(1),\qquad
 \lambda_{t,s}(A)=[A]_2\prod_{j=s+1}^{t-1}[p_j]_1^2\sqrt{k_t/k_s}.
\end{equation}
从而 $J_{t,N}[A;p]$ 等于 $R_{t,N}[A;p]$ 加上有限个深度严格小于 $t$
的原始正次多项式族与 $o_{L^2}(1)$；每个较浅族只改变根多项式。
所有确定性系数在固定宽比下有界。
\end{lemma}
\begin{proof}
按从高到低的顺序逐一恢复主族的界面掩码，暂不恢复产生的修正族。
在界面 $s$，高于它的掩码已经全部移除，故可对上部全原始系数使用
\eqref{fp:actual-raw-interface} 与\eqref{fp:actual-diag}。
一次项无碰撞，高阶碰撞为 $o(1)$；二阶主对角代入
$y_i=p_s(z_i)$，恰得
$\lambda_{t,s}(A)k_s^{-1/2}\sum_i p_s(z_i)^2$。
低于 $s$ 的所有掩码仍保留，所以这正是式中带负号的
$J_{s,N}[p_s^2;p_{<s}]$，两个共同标签因子先精确相乘。
如此有限次恢复即得\eqref{fp:actual-triangular}。
误差经过的完整提升范数至多为固定个 $\|A_j\|_{\rm op}$ 的固定次幂，
掩码和对称化都是收缩；Gaussian 固定次数矩等价、H\"older 不等式与
式~\eqref{fp:actual-isometry} 保持 $o_{L^2}(1)$。
再对严格较浅的每个修正族归纳，最后到 $J_{1,N}=R_{1,N}$，即得后一结论。
全过程只对全原始上部核使用碰撞定理。
\end{proof}

\paragraph{同一二次系数的精确消去。}
式~\eqref{fp:actual-triangular}中的系数满足
\[
 \lambda_{t,s}(A)=\lambda_{t,t-1}(A)\lambda_{t-1,s}(p_{t-1}^2)
 \quad(s<t-1),\qquad
 [p_{t-1}^2]_2=[p_{t-1}]_1^2.
\]
把深度 $t-1$、根为 $p_{t-1}^2$ 的同一三角恒等式代入，即得
\begin{equation}\label{fp:actual-restoration-exact}
 J_{t,N}[A;p]=R_{t,N}[A;p]
 -[A]_2\sqrt{k_t/k_{t-1}}\,
 R_{t-1,N}[p_{t-1}^2;p_{<t-1}]+o_{L^2}(1).
\end{equation}
这是确定主对角的精确代数消去；有限维误差仍按前述完整提升估计处理。

对至多多项式增长的 $f$，Gaussian 密度关于方差求导给
$\frac{d}{dq}\mathbb E f(\sqrt qG)|_{q=1}
=\frac12\mathbb E[f(G)(G^2-1)]$，所以
$[H_f]_2=\lambda_f$（$f=\sigma,\chi$）。
因而顶根递推和对角恢复合成同一个局部关系
\begin{equation}\label{fp:actual-cancel-recursion}
 \Pi_NF_t^f=R_{t,N}[H_f;P_N]
 +\lambda_f\sqrt{k_t/k_{t-1}}
 \bigl(\Pi_NF_{t-1}^{\chi}-R_{t-1,N}[P_N^2;P_N]\bigr)+o_{L^2}(1).
\end{equation}
置 $U_N=[Q_N-P_N^2]_{\deg\le N}$，
$E_t=\Pi_NF_t^\chi-R_{t,N}[P_N^2;P_N]$。
由 $\lambda_\chi=1$ 和深度一结论，
\[
 E_1=R_{1,N}[U_N]+o_{L^2}(1),\qquad
 E_t=R_{t,N}[U_N;P_N]+\sqrt{k_t/k_{t-1}}E_{t-1}+o_{L^2}(1).
\]
固定深度归纳，再将 $f=\sigma$ 代入\eqref{fp:actual-cancel-recursion}，得到
\begin{equation}\label{fp:actual-U}
 \mathcal Q_N=R_{L,N}[P_N;P_N]
 +\mu'\sum_{t=1}^{L-1}\sqrt{k_L/k_t}\,R_{t,N}[U_N;P_N],
 \qquad \mu'=bu/2,
\end{equation}
以及
\begin{equation}\label{fp:actual-final}
 \boxed{\displaystyle
 \mathbb E_W\|\Pi_NF_L-\mathcal Q_N\|_{L^2_X}^2\longrightarrow0,
 \qquad N>2.}
\end{equation}
奇分支条件均值为零，且原始界面次数为奇数而无二阶对角，故直接取
$\mu'=0$，不需引入 $\chi$ 或 $E_t$。
所有根多项式没有常数项；正次代入不能降低形式次数，故可统一截到 $N$。

\paragraph{任意先验边缘耦合。}
若 $(W,W')$ 的两个边缘均为网络 Gaussian 先验，且
$\sup_d\mathbb E\|B(W)\|_2^2<\infty$，则
\[
 \mathbb E|\langle B(W),\Pi_NF_{W'}-\mathcal Q_N(W')\rangle|
 \le(\mathbb E\|B(W)\|_2^2)^{1/2}
      (\mathbb E\|\Pi_NF_{W'}-\mathcal Q_N(W')\|_2^2)^{1/2}=o(1).
\]
这里只用各自边缘的误差，不在耦合下额外假设层间独立。
闭合保留全部主深度原始 profile 和有限较浅根族；同深度根合并后再作全剖面比较。

\section{有序树、物理槽分配与全剖面分离}\label{fp:history}\label{app:cuts}
\newthought*{固定顺序再求和。}
本节固定深度 $L$、次数 $N=m^L$ 和比例宽度 $k_s\asymp d$。奇激活取奇数 $m\ge5$；分段线性激活取偶数 $m\ge6$。所有常数允许依赖这些固定参数。用 $K_d$ 表示一个固定的对数幂，令 $\delta_d=K_d/\sqrt d$、$\kappa_d=K_d/\sqrt d$。以下所有有限次数、高概率界均在源、教师各自的原始高斯边缘事件上成立；两个网络之间没有独立性假设。

一个\emph{有序剖面}是在每个原始计算树节点选择一个正次数单项式，并给该节点的孩子指定顺序。物理槽是完全展开后按顺序排列的 $N$ 个输入向量位置。固定一个把这些位置分配到平衡教师 $N$ 个位置的双射；所有估计先对这个固定双射证明。张量对称化只是至多 $N!$ 项的固定有限平均，不能先以对称化掩盖某项中的相干收缩。不同节点出现可以选择不同的允许次数，重复的神经元数值标签与不同的树节点出现也必须区分。

第 $h$ 相的物理槽表示已经完成的高度 $h-1$ 教师块；第 $1$ 相的槽就是输入位置。每个待完成教师叶含 $m$ 个当前槽。完整收缩这 $m$ 个槽后，源节点上得到一个\emph{字面因子} $M_{ia}$：它保留源神经元标签 $a$ 和一个完成教师块标签 $i$，以及全部原始边的归一化。它不是一条新独立高斯边。同相刚产生的字面因子只能作为旁观输出，下一相才可成为活动槽。

\subsection{教师上下文与最低切口}
\newthought*{保留完整余张量。}
教师当前原始齐次系数记作 $C$。完成的输出位置采用恒等映射，未完成位置采用有界齐次幂列映射。固定一个活动教师行 $i$ 后，其余全部物理输出的范数因此至多为 $K_d\|C(e_i,\cdot)\|_{\rm HS}$，而
\begin{equation}\label{fp:teacher-slices}
 \sum_i\|C(e_i,\cdot)\|_{\rm HS}
 \le \sqrt{k}\|C\|_{\rm HS}\le K_d\sqrt d.
\end{equation}
交叉度量始终放在源侧；已经自适应收缩的教师张量不再被声称是新的原始齐次系数。

在源切口的每个停止节点，把其严格下方的原始子树和已安装输出一同作为一个完整有界特征映射。严格高于停止高度的每个字面因子单独保留一个系数输出模式。即使多个输出共享同一源标签，也保留逐输出的恒等复制模式。结构复制不是不同源父节点，不能删除其对角。当前高度的所有节点都必须在特征映射和恒等复制\emph{之前}停止，才能使用新颜色的右旋转估计。

\begin{lemma}[最低非平凡相遇]\label{fp:lowest-cut}
假定旧字面因子在第一次活动使用之前，已对带有剩余因子的情形作完整余张量检验。若所有含两个活动源父位置的切口都已被相应解析估计排除，则每个存活的新教师叶收缩，恰好在一个原始源父节点消耗 $m$ 个完整、在当前变量中线性的原始孩子因子。
\end{lemma}
\begin{proof}
给待完成教师叶的 $m$ 个源槽逐一作标记。按原始源高度递增推动这些标记，起点较高的标记暂时等待。标记在遇到第一个乘法分支之前，经过的表达式是原始归一化矩阵的连续乘积，末端接其已有度量，故仍是有界行映射。删除一个选中线性因子后若没有物理因子或旁观输出，该父节点称为单因子节点。到达同一父节点的标记首先分组；两个标记同时到达的节点必然不是单因子节点。

若遇到的节点都是单因子节点，则精确求和其标签并沿原始线性路径继续。否则在最小高度的非单因子节点停止。到达这里的每个选中因子都是完整线性因子：若在其下方已穿过非线性分支，则应更早停止。若标记分布在至少两个父位置，选其中最低的非单因子位置作为新颜色坐标。较低的另一单路径推进到该高度；若它更早遇到分支，则改以该更早分支停止。较高活动旧字面因子已经通过首次使用检验，因而不带剩余因子，不可能还是另一活动标记的祖先；本相新字面因子只是旁观输出。由原始树重新构造共同切口，并按上文规则在本高度所有特征映射之前停止，便获得所需的新颜色坐标切口。

按假设，该混合父位置的情况被排除。因此全部标记在同一父节点相遇。如果其中两个标记来自同一个非线性原始孩子，它们在严格更低节点就已相遇并应停止，矛盾。所以它们来自 $m$ 个不同的完整原始孩子，且每个孩子在当前物理变量中线性。
\end{proof}

切口阶数至少三时，最低非单因子位置是 fresh 模式，删除一个选中线性槽后
仍有非空残余，因此直接适用双父收缩估计~\ref{fp:core-two-parent}。
所有配对都在恢复后的原始 profile 上进行，包含相同数值标签，
不再在本步骤插入碰撞掩码。每项完整余张量界为 $K_d/d$，与
式~\eqref{fp:teacher-slices} 配对的误差为 $K_d/\sqrt d=o(1)$。

\subsection{奇数剖面的闭合历史}
\newthought*{奇偶性用于切口阶数。}
奇激活中所有原始次数为奇数。每个原始替换将输出数改变偶数 $r-1$，一个已证明完整的新出生将 $m$ 因子替为一个，改变偶数 $-(m-1)$。因此每个非空停止子树的当前物理与旁观输出总数为奇数，把它停止为一个模式也保留奇偶性。此归纳仅使用\emph{以前}已接受的出生，绝不假定正在检验的出生完整。于是有至少两个不同活动父位置的展开切口，阶数必为至少三。严格较高的每个字面输出仍单独计数，结构复制只计输出模式、不计独立活动父位置。

原始次数为 $r$ 的节点若本相已发生 $q$ 次出生，因子列表含 $q$ 个字面因子和 $r-qm$ 个未消耗原始孩子，可见因子数为 $r-q(m-1)$，仍为奇数。旧字面因子第一次活动使用时若尚带剩余因子，则至少还有两个因子。若本次教师叶还访问另一源列，引理~\ref{fp:core-split-hit} 的分裂命中估计排除该项；若全部命中本列，当上方源系数阶数至少二时使用引理~\ref{fp:core-hybrid-Rge2}，当阶数为一时使用引理~\ref{fp:core-hybrid-R1} 的至少两个额外因子估计。额外字面因子也分别保留；普通非线性因子保留原始高斯读出结构。因此存活的活动旧字面因子总是其出生节点的全部特征。这个结论提供引理\ref{fp:lowest-cut}的首次使用前提。

完整出生的度量是 $m$ 个线性行核的逐项乘积。若核为 $\rho^{(j)}$，由各行各列的有界 $\ell^2$ 范数和 H\"older 不等式，
\begin{equation}\label{fp:metric-schur}
 \sup_i\sum_a\left|\prod_{j=1}^m\rho^{(j)}_{ia}\right|
 +\sup_a\sum_i\left|\prod_{j=1}^m\rho^{(j)}_{ia}\right|\le K_d.
\end{equation}
故度量的绝对行列和有界。下一教师原始高斯矩阵 $A$ 满足 $\max|A_{ij}|\le\kappa_d$，所以即使度量依赖两网络，仍有 $\max|AM|\le K_d\kappa_d$，无需条件独立性。

每次路径推动暴露有限原始树中的一条边，每次出生完成一个教师块，因此每相有限终止。一相结束后以上述度量作下一相的物理终端，先进行首次使用检验，再继续最低切口。最终出生不可能带正次数剩余因子，因为全部 $N$ 槽已经用完。

高度 $h$ 的存活块需要沿每条路径至少 $h$ 个原始源节点：$h=1$ 需要分支节点，$h$ 的 $m$ 个完整孩子分别是高度 $h-1$ 的块，而活动使用时无剩余因子，故父出生严格位于这些孩子之上。任何一元延迟只会增加源高度。深度至多 $L$ 的源若产生高度 $L$ 教师块，则每条路径都取等号，没有一元延迟，所有节点次数恰为 $m$。因此奇数剖面中只有完整平衡 $m$ 元剖面存活；其表达式保持全部原始系数与归一化，正是齐次比较项。

\subsection{永久标记路径与二模式的单因子边界}\label{fp:permanent-mark}
\newthought*{不提前收缩标记路径。}
考虑一个精确二模式状态
\begin{equation}\label{fp:marked-state}
 \sum_{a,b}E_{ab}\,\ell_a\otimes Q_b.
\end{equation}
其中 $\ell$ 是一个完整当前教师终端因子，$Q$ 是原始高度至多 $L-1$ 的完整子树；所有尚未使用的物理和旁观因子均保留。$E$ 必须来自原始二次切口或其原始矩形传输，而非任意自适应矩阵；其原始重构系数满足引理~\ref{fp:core-closure} 的核心估计。特别地 $\|E\|_{\rm HS}\le K_d$，各非平凡展平至多 $\delta_d$。一个较短源根留下一个字面因子与一个完整孩子时，就是 $E_{ab}=d_aU_{ab}$，其中 $\|d\|_2\le K_d$、$\|d\|_\infty\le\kappa_d$，$U$ 为原始归一化高斯边。

永久保留 $a$ 模式，标记从 $\ell$ 到最终教师根的唯一路径，永不关闭该路径上的教师块。其余教师块按高度递增处理，它们所有源槽都在 $Q$ 内。每次重构核心仍保留 $a$，且当前新颜色停止在全部复制和特征映射之前。于是选中一个后代位置时核心阶数至少二，选中两个不同后代位置时至少三。

因此 $Q$ 内的首次使用检验总可使用源展平版本；固定教师行的完整源余张量范数至多 $K_d\delta_d(\delta_d+\kappa_d)$。所有标记路径的剩余槽包含在该余张量中，乘\eqref{fp:teacher-slices}后趋零。混合父位置检验则用双父收缩估计~\ref{fp:core-two-parent}。存活的非标记出生仍消耗完整孩子。它不能发生在 $a$：其全部活动槽在唯一孩子 $Q$ 内，若该孩子非线性则更早在内部相遇，若它只有一个当前槽则不足以提供 $m$ 槽。

每个非标记高度 $h$ 出生需要 $Q$ 内原始高度至少 $h$。若 $Q$ 高度小于 $L-1$，最终根的 $m-1$ 个非标记高度 $L-1$ 孩子不可能完成。若 $Q$ 高度等于 $L-1$，这些 $m-1$ 个孩子必须全部出生在 $Q$ 的唯一原始根 $b$。令其完整度量为 $M^{(j)}$，最终教师矩阵为 $A$，并置 $G^{(j)}=AM^{(j)}$。由\eqref{fp:metric-schur}，每个 $G^{(j)}$ 的列 $\ell^2$ 范数有界，最大条目至多 $K_d\kappa_d$。

保持标记教师孩子完整物理化。固定最终教师行 $i$，它的范数有界。去掉 $m-1$ 个已知字面因子后，完整源余张量为
\[
 R_b=\sum_aE_{ab}\ell_a\otimes Q_b^{\rm res},\qquad
 \|R_b\|\le K_d\|E_{:b}\|_2,
 \quad\sum_b\|R_b\|\le K_d\sqrt d\|E\|_{\rm HS}.
\]
这里每个余因子仍原样保留，$\ell$ 的列映射有界，$Q_b^{\rm res}$ 的列范数有界。归一化最终读出最大值为 $O(d^{-1/2})$；利用两个 $G$ 的列 $\ell^2$ 界及其余 $m-3$ 个的条目界，完整标量配对至多
\begin{equation}\label{fp:marked-terminal}
 K_d\|E\|_{\rm HS}\kappa_d^{m-3}=o(1).
\end{equation}
这一步同时使用 $m-1$ 个字面因子，未把一个小标量核与任意余张量直接相乘。它解决了完整剩余深度下的单因子/非单因子二模式边界，也解决较短任意根的一个字面因子加一个完整孩子的边界。

\subsection{分段线性剖面与相干二次分支}
\newthought*{对源高度作强归纳。}
普通多项式 $P$ 只含一次与正偶次。固定偶数 $m\ge6$。同时证明：每个高度 $t<L$、根为任意正次数多项式 $A$、其下各节点为 $P$ 的任意有序单项式选择，在每个固定槽分配下配对趋零；高度 $L$ 的普通 $P$ 剖面只有全 $m$ 剖面存活。归纳允许不同孩子选不同单项式，以覆盖融合后的精确有序剖面。

最低切口阶数至少三已处理；阶数一不能有两个不同活动父位置。阶数二时，一单因子、一非单因子的情况由\eqref{fp:marked-state}--\eqref{fp:marked-terminal}解决。若两者都非单因子，则它们必须都是同高度的新模式：较高活动字面因子首次使用后已无剩余因子，而较高旁观字面输出会成为第三模式。该切口故是尚未施加下面特征映射的原始二次上方系数 $D$。

对它的非对角部分，由置换对称性、有界 Hilbert 二阶矩和固定次数高斯超可积性，有
\[
 \max_{a\ne b}|D_{ab}|\le K_d/d.
\]
保留两标签后，完整源余张量界为
\[
 K_d\Bigl(\sum_{a\ne b}|\eta_{ia}\zeta_{ib}D_{ab}|^2\Bigr)^{1/2}
 \le K_d/d,
\]
与教师切片和配对即趋零。该论证只使用坐标/坐标估计，未使用一般二阶坐标/向量增益。

对于对角，式~\eqref{fp:actual-diag} 的原始二次系数估计给出
\begin{equation}\label{fp:coherent-diagonal}
 \operatorname{Diag}D=\lambda_s I/\sqrt{k_s}+E_s,
 \qquad\mathbb E\|E_s\|_{\rm HS}^2=o(1).
\end{equation}
先通过完整有界源特征映射支付误差。主项识别两个原始父标签，把它们的两个选中根单项式相乘，保留所有有序孩子与完全相同的教师槽分配。得到高度 $s<t$、读出 $1/\sqrt{k_s}$ 的原始前缀剖面，乘精确确定性系数 $\lambda_s$。其根可有任意奇偶性，所以必须使用上述加强归纳类。已完成的字面度量可展开回原始槽收缩；它们是等式，故融合不产生残缺子树。若 $s=0$，只剩输入次数二，不能贡献 $N>2$。

二模式也可能含一个较高字面输出与一个新模式：它没有同色对角删除，活动情况下属于永久标记论证，旁观情况下只有一个活动父模式。如果两模式都是较高纯字面输出，则未完成状态中每个最多代表 $m^{L-1}$ 个原始槽，两者总数小于 $N$，不能是整个剖面；若一个已为高度 $L$，它独自已用尽全部槽。结构复制同样按其代表的原始槽数计算，不能当成新的独立父节点。

普通偶次节点在本相 $q$ 次出生后剩余 $r-qm$ 个原始孩子，该数为偶数。一个活动字面因子若带正剩余，就至少还有两个普通因子；或有其他分别保留的字面因子。因此上方核心阶数一时乘积余张量估计适用，阶数至少二时源展平估计适用。纯字面因子列同样服从上述三项：若本叶还命中其他列，使用第一项；
否则本叶的 $m$ 个当前物理槽全部来自该列，该列至少有 $m$ 个单槽因子，
所以一模时有至少两个额外因子，适用第三项，多模时适用第二项。
特别地，仅有两个字面因子的列不能供给一个 $m\ge6$ 槽教师叶，必属第一项。
任意较短根不要求次数偶性，其首次使用后的因子列表写成 $q$ 个字面因子与 $k=r-qm\ge0$ 个原始孩子：若 $q+k\ge3$ 则至少两个额外因子；若 $q+k=2$，只可能两个字面因子，或一个字面因子加一个完整原始孩子，后者由永久标记论证处理；若 $q+k=1$，则为无剩余的完整字面因子。这穷尽全部根次数。

从 $t=1$ 开始归纳，所有二次对角只落入严格更低高度；没有正高度的更低项时，输入次数二排除该项。排除混合切口、首次使用的剩余因子后，有限调度使每个出生完整且活动时无剩余。高度 $t<L$ 不能产生高度 $L$ 的最终块，故其配对趋零。随后处理高度 $L$ 的主剖面，所有融合项都已在较短归纳中排除。高度取等号迫使全树无一元延迟且每个节点次数为 $m$；其存活项恰为原始齐次项。

\subsection{从确定性事件到任意先验耦合}\label{fp:coupling-uniform}
\newthought*{边缘误差先支付。}
固定次数和深度时，有序剖面、槽分配、原始切口、融合高度和历史形状的数量均与 $d$ 无关。分别在两个边缘上同时施加所有前述原始高斯事件，然后才允许交叉度量依赖两个网络。在共同好事件上，各项估计均为确定性 $o(1)$。在坏事件上，用固定次数系数的有界边缘高阶矩和 H\"older 不等式；坏事件概率可取任意所需多项式阶。该控制不依赖耦合，因此对所有先验边缘耦合一致。

设 $T_W$ 为与激活对应、带完整标量系数的全 $m$ 平衡齐次输入混沌项，并记式~\eqref{fp:actual-U} 的有限原始多项式家族为 $\mathcal Q_N(V)$。上文逐剖面结论给出
\begin{equation}\label{fp:profile-pairing}
 \sup_{\pi\in\Gamma(\mathsf P_W,\mathsf P_V)}
 \mathbb E_\pi\left|\langle T_W,\mathcal Q_N(V)-T_V\rangle\right|
 \longrightarrow0.
\end{equation}
又 $T_W$ 完全位于第 $N$ 输入混沌，故 $\langle T_W,F_V\rangle=\langle T_W,\Pi_NF_V\rangle$。由式~\eqref{fp:actual-final}，自身边缘的闭合误差用 Cauchy--Schwarz 支付：
\[
 \mathbb E_\pi|\langle T_W,\Pi_NF_V-\mathcal Q_N(V)\rangle|
 \le(\mathbb E\|T_W\|^2)^{1/2}
      (\mathbb E\|\Pi_NF_V-\mathcal Q_N(V)\|^2)^{1/2}=o(1).
\]
结合\eqref{fp:profile-pairing}得到所需全局函数比较
\begin{equation}\label{fp:global-function-pairing}
 \boxed{\displaystyle
 \sup_{\pi\in\Gamma(\mathsf P_W,\mathsf P_V)}
 \mathbb E_\pi|\langle T_W,F_V-T_V\rangle|\longrightarrow0.}
\end{equation}

\section{齐次比较网络的全局匹配}\label{hm:section}\label{app:matching}
\newthought*{本节的对象与结论}
本节仅讨论激活为 $c_m z^m$、$c_m=(m!)^{-1/2}$ 的齐次网络。固定 $L$ 与整数 $m\ge5$，令 $k_0=d$，$cd\le k_s\le Cd$，并定义
\[
 h^{(0)}=x,\qquad h^{(s)}=c_m(W^{(s)}h^{(s-1)}/\sqrt{k_{s-1}})^{\odot m},\qquad
 \mathcal H_W(x)=k_L^{-1/2}\sum_i h_i^{(L)}(x),\qquad
 p=\sum_{s=1}^L k_{s-1}k_s.
\]
每个网络的所有权重分别具有独立标准高斯先验；两个网络之间可任意耦合。记 $P_s(Y)$ 为将 $h^{(s)}$ 替换为形式变量 $Y$ 后的下游输出，$q_s=m^{L-s}$，$C_s=C_{q_s}(P_s)$。这里 $P(Y)=\langle C_r(P),Y^{\otimes r}\rangle$，张量使用继承的 Hilbert 范数；若 $P=\sum_{|\nu|=r}p_\nu Y^\nu$，则
\begin{equation}\label{hm:coeff-norm}
 \|C_r(P)\|^2=\sum_{|\nu|=r}\frac{\nu!}{r!}|p_\nu|^2.
\end{equation}
特别地 $C_L=k_L^{-1/2}{\bf1}$。

\begin{theorem}[齐次网络的有限标签全局匹配]\label{hm:global-matching}
存在一族确定性的参数对齐映射 $\mathcal G_d$，以及仅依赖固定深度、次数和宽度比例的常数 $K$，使得
\begin{equation}\label{hm:main}
 \log|\mathcal G_d|=O(d\log d)=o(p),\qquad
 \frac1p\mathbb E\min_{g\in\mathcal G_d}\|W-gV\|_F^2
 \le K\mathbb E\|C_0(W)-C_0(V)\|^2+o_d(1).
\end{equation}
余项对所有具有上述先验边缘分布的耦合一致。若 $X\sim N(0,I_d)$，则右侧亦可替换为
$K\mathbb E\|\mathcal H_W-\mathcal H_V\|_{L^2_X}^2+o_d(1)$。
\end{theorem}

\subsection{所需共同框架及下游模态能量}
\newthought*{精确列出共同框架的用法}
本 PDF 的重数框架定理（框架界~\eqref{tg:frame}）在本节只以如下形式使用。对单位行向量 $a_i$、$k=O(d)$，若最大自身相干度 $\eta=O(\sqrt{\log d/d})$，固定 $r$，则纯张量
\[
 \operatorname{Sym}\bigotimes_i a_i^{\otimes m\nu_i},\quad |\nu|=r,
\]
以及单删除张量
\[
 \operatorname{Sym}\bigotimes_i a_i^{\otimes(m\nu_i-\mathbf1_{i=j})},
 \quad |\nu|=r,\quad \nu_j>0,
\]
各自构成具有常数上下界的框架（不同的 $(\nu,j)$ 无重复）。即其任意线性组合的范数平方与系数平方和在维数无关常数内等价。容许 $1/2\le\|a_i\|\le2$，只改变常数。充分的小扰动界为
$K[\eta+\sum_{h=1}^r k^h\eta^{mh-2}]=o(1)$。
这里全局对称化不能换成带颜色的张量空间。以下所有其他比较均在本节证明。

令 $B_{s,i}=\partial_iP_s$。在各层行范数、相干度和列平方和的通常高斯集中事件上，
\begin{equation}\label{hm:derivative}
 a/k_s\le\|C(B_{s,i})\|^2\le b/k_s\quad\hbox{对所有 }s,i.
\end{equation}
为完整说明这一共同框架的推论，终层为 $B_{L,i}=k_L^{-1/2}$。置 $w_j=W^{(s+1)}_{j:}/\sqrt{k_s}$，链式法则给出
\[
 B_{s,i}(Y)=mc_m\sum_j w_{ji}(w_j\cdot Y)^{m-1}
 B_{s+1,j}\big(c_m(w\cdot Y)^{\odot m}\big).
\]
展开 $B_{s+1,j}$ 后，每项恰为标签 $(\nu,j)$ 的单删除张量；固定次数下式 \eqref{hm:coeff-norm} 与普通单项式系数范数等价。删除框架因此给出
\[
 \|C(B_{s,i})\|^2\asymp\sum_j w_{ji}^2\|C(B_{s+1,j})\|^2
 \asymp k_{s+1}^{-1}\sum_j w_{ji}^2\asymp k_s^{-1}.
\]
这是确定性框架估计，无须下游导数与当前行独立。若 $H_s$ 为 $C_s$ 的第一模态 Gram 矩阵，则精确微分恒等式给出
\begin{equation}\label{hm:mode}
 (H_s)_{ii}=q_s^{-2}\|C(\partial_iP_s)\|^2,\qquad
 a'/k_s\le(H_s)_{ii}\le b'/k_s.
\end{equation}

\subsection{稠密系数与坐标碰撞的控制}
\newthought*{使自适应删除合法的边缘事件}
对任意行矩阵 $A$，置 $U_Ae_i=a_i^{\otimes m}$。由于
$U_A=(A^T)^{\otimes m}J$，其中 $Je_i=e_i^{\otimes m}$ 为等距嵌入，
$\|U_A\|_{\rm op}\le\|A\|_{\rm op}^m$。每次齐次代入为 $U_A$ 的张量幂再接收缩映射 $\operatorname{Sym}$。从 $\|C_L\|=1$ 递推，$\|C_s\|$ 被有限个归一化高斯矩阵算子范数的固定幂乘积控制。高斯矩阵算子范数的固定阶矩有一致上界，故
\begin{equation}\label{hm:moments}
 \sup_d\mathbb E\|C_s\|^u<\infty\quad(0<u<\infty).
\end{equation}
对平方自由单项式 $Y_I=\prod_{i\in I}Y_i$，$|I|=q_s$，写其系数为 $p_I$，置 $c_I=p_I/\sqrt{q_s!}$。式 \eqref{hm:coeff-norm} 表明 $\sum_Ic_I^2$ 为该扇区的张量范数平方。上层第一矩阵的列置换不变性及 \eqref{hm:moments} 给出
$\mathbb E c_I^2\le C\binom{k_s}{q_s}^{-1}\le Ck_s^{-q_s}$。
每个 $c_I$ 是次数不超过 $D_s=m+\cdots+m^{L-s}$ 的高斯多项式，故高斯多项式超收缩不等式
$\|c_I\|_u\le(u-1)^{D_s/2}\|c_I\|_2$（$u\ge2$）给出：取 $u=A\log d$，在 $e\|c_I\|_u$ 处应用 Markov，再对至多 $k_s^{q_s}$ 项作并集界，可令任意预定 $M$ 下
\begin{equation}\label{hm:dense}
 \mathbb P\left\{\max_I|c_I|^2>
 C_M (\log d)^{D_s}k_s^{-q_s}\right\}\le d^{-M}.
\end{equation}
终层为确定向量，单独处理。

还需 \emph{整个} 下游张量与平方自由截断相差 $o(1)$。以下给出直接算子证明。令 $T=\operatorname{Sym}_{mr}U_A^{\otimes r}B$，$\|A\|_{\rm op}\le K$，$\max_{ij}|a_{ij}|\le\kappa$。将 $T$ 的两个指定输入槽与 $e_j$ 收缩：若两个槽来自同一个 $m$ 块，所得块映射为
\[
 (A^T)^{\otimes(m-2)}J\operatorname{diag}(a_{ij}^2),
\]
其范数至多 $K^{m-2}\kappa^2$；若来自不同块，每块各自范数至多 $K^{m-1}\kappa$。其余块的算子范数有常数上界，固定次数下对称化只有有限种槽分配。因此
\[
 \|T(e_j,e_j,\cdot,\ldots,\cdot)\|\le C\kappa^2\|B\|.
\]
坐标碰撞张量范数平方至多所有坐标对角切片范数平方之和乘以 $\binom{mr}{2}$，所以
\begin{equation}\label{hm:collision}
 \|(I-\mathrm{SF})T\|^2\le C n\kappa^4\|B\|^2=o(1)\|B\|^2,
 \qquad \kappa=O(\sqrt{\log d/d}),\ n\asymp d.
\end{equation}
将此用于 $C_s$ 的第一次下游代入，并用上层张量范数有界，即得
$\|(I-\mathrm{SF})C_s\|=o(1)$。不能由此推出平方自由截断的逐坐标模态下界；后文先恢复完整张量再用 \eqref{hm:mode}。

所有这些边缘事件可同时选取，失败概率趋于零。在行有符号置换和列有符号置换下，行范数、算子范数、相干度、最大坐标、列平方和不变；下游张量仅重排并改变坐标符号，故范数、系数绝对最大值、模态对角上下界、碰撞扇区也不变。于是自适应选择内部对齐后这些事件仍成立；无需断言变换后的网络仍为独立高斯。

\subsection{全局对称与分组张量的精确换算}
\newthought*{半径与阶乘不能遗漏}
在第 $s$ 层写 $w_i=W^{(s)}_{i:}/\sqrt{k_{s-1}}=r_i a_i$，$\|a_i\|=1$，
$R=\operatorname{diag}(r_i^m)$，$\overline C_s=R^{\otimes q_s}C_s$。精确代入为
\begin{equation}\label{hm:substitution}
 C_{s-1}=c_m^{q_s}\operatorname{Sym}_{mq_s}U_A^{\otimes q_s}\overline C_s.
\end{equation}
均匀半径误差 $\max_i|r_i-1|\le\delta_d\to0$ 意味着
$\|\overline C_s-C_s\|\le C\delta_d\|C_s\|$。
模态能量、最大系数与碰撞估计仍成立，因为每个张量坐标只乘上固定个 $1+o(1)$ 半径因子。

固定 $q=q_s$，对 $|I|=q$ 定义
\[
 Q_I=\sqrt{q!}\operatorname{Sym}_q\bigotimes_{i\in I}a_i^{\otimes m},\qquad
 E_I=\sqrt{(mq)!/(m!)^q}\operatorname{Sym}_{mq}\bigotimes_{i\in I}a_i^{\otimes m}.
\]
$Q_I$ 属于分组空间，$E_I$ 属于实际全局对称空间。平方自由坐标基向量经分组提升变为 $Q_I$，再经 \eqref{hm:substitution} 变为 $\sqrt{\lambda_{m,q}}E_I$，其中
\begin{equation}\label{hm:lambda}
 \lambda_{m,q}=q!/(mq)!.
\end{equation}
两种归一化使不拆分 $m$ 个槽的跨网络配对项完全相同。

\subsection{自适应混合槽比较}
\newthought*{一次 Schur 比较同时处理自身与交叉 Gram。}
给定两个单位行阵列 $A,B$，自身相干度至多 $\eta_d=O(\sqrt{\log d/d})$，
算子范数至多 $K$。置 $\rho_{ij}=\langle a_i,b_j\rangle$、$\tau=d^{-1/8}$。
二部图 $|\rho_{ij}|>\tau$ 的度数至多 $D=\lceil K^2/\tau^2\rceil$；
在各侧连接共享邻居的顶点，所得图度数至多 $D^2$。不含图边的 $q$ 元集合称为好集合。
对好 $I=\{i_1,\ldots,i_q\}$，任一列中至多一个 $|\rho_{i_a j}|>\tau$，故
\begin{equation}\label{hm:mixed-column}
 \sum_j\prod_a|\rho_{i_a j}|^{v_a}\le
 \begin{cases}K^2,&\#\{a:v_a>0\}=1,\\ C\tau,&\#\{a:v_a>0\}\ge2,
 \end{cases}\qquad \sum_a v_a=m.
\end{equation}
后一界按最大正支撑相关系数的位置分区：其指数至少二时保留该项的平方，
否则保留它与另一项各一幂；剩余因子中至少一幂至多 $\tau$。
所保留项的列和分别由平方和界及 Cauchy--Schwarz 控制。
全槽配对计数给出
\[
 \langle E_I^A,E_J^B\rangle
 =\sum_{\substack{e_{ab}\ge0\\\sum_b e_{ab}=\sum_a e_{ab}=m}}
 \frac{(m!)^q}{\prod_{a,b}e_{ab}!}\prod_{a,b}\rho_{i_a j_b}^{e_{ab}}.
\]
其中轮廓 $e=mP$ 恰给 $\langle Q_I^A,Q_J^B\rangle$；其他轮廓至少有一混合列。
固定好 $I$，把对好 $J$ 的绝对行和放大为对所有有序 $J$ 的和，各列分离，
由 \eqref{hm:mixed-column} 得 $C\tau$；交换两侧得到同一列和界。因此
\begin{equation}\label{hm:cross-schur}
 \|K_{\rm act}-K_{\rm grp}\|_{\rm op}\le C\tau
\end{equation}
在双方好集合上成立。这也适用于 $A=A$、$B=B$；因 $\eta_d<\tau$，这两种自身比较的所有集合均好。
最大度数 $D^2$ 的图中，含边 $q$ 元集合至多 $C_qD^2k^{q-1}$ 个。
由 \eqref{hm:dense}，同时对所有可能依赖完整网络及系数的图，删除能量至多
$C D^2k^{q-1}\max_I|c_I|^2\le C d^{-1/2}(\log d)^{D_s}=o(1)$。
实际与分组提升的算子范数均有界，故删除以及 \eqref{hm:collision} 的平方自由截断
各只产生 $o(1)$ 范数误差。对三个 Gram 块 $AA,AB,BB$ 使用同一比较，再恢复完整系数，得
\begin{equation}\label{hm:sym-group}
 \|C_{s-1}^A-C_{s-1}^B\|^2
 =\lambda_{m,q}\|G_A-G_B\|^2+o(1),\qquad
 G_A=U_A^{\otimes q}\overline C_s^A,\quad
 G_B=U_B^{\otimes q}\overline C_s^B.
\end{equation}
误差在联合典型事件上一致；$q=1$ 时无拆槽，此等式精确成立。

\subsection{一个平方差恒等式同时给出匹配与系数恢复}
\newthought*{绝对行列和控制替代投影与提升扰动。}
记 $U=U_A,V=U_B,M=U^*V=(\rho_{ij}^m)$，$\rho_i=\max_j|\rho_{ij}|$。
一行中 $|\rho_{ij}|>\tau$ 的子集大小至多 $D$，其自身 Gram 算子范数至多 $1+D\eta_d$；
低相关项保留平方后提出 $\tau^{m-2}$，故
\begin{equation}\label{hm:cross-power}
 \sum_j|M_{ij}|\le\rho_i^{m-2}(1+D\eta_d)+K^2\tau^{m-2}
 \le\rho_i^{m-2}+e_d,\qquad \sum_i|M_{ij}|\le1+e_d,
\end{equation}
其中 $e_d\to0$ 对两阵列一致；列界由同一论证得到。
置 $\widetilde M=M/(1+e_d)$、$\mathsf K=\widetilde M^{\otimes q}$，
并将完整系数张量 $C=\overline C_s^A,D'=\overline C_s^B$ 写为有序指标数组 $(c_I),(d_J)$。
$\mathsf K$ 的绝对行列和 $r_I,s_J$ 均不超过一，且有精确恒等式
\begin{equation}\label{hm:transport-identity}
 \begin{aligned}
 \mathcal E&:=\|C\|^2+\|D'\|^2-2\langle C,\mathsf K D'\rangle\\
 &=\sum_I(1-r_I)c_I^2+\sum_J(1-s_J)d_J^2
   +\sum_{I,J}|\mathsf K_{IJ}|
       (c_I-\operatorname{sgn}(\mathsf K_{IJ})d_J)^2.
 \end{aligned}
\end{equation}
零权项不影响等式。因为 $U^*U,V^*V=I+o_{\rm op}(1)$（非对角绝对行和至多 $k\eta_d^m=o(1)$），
且系数范数有界，$\mathcal E=\|G_A-G_B\|^2+o(1)$。
$r_I$ 是各槽行和的乘积，故 $1-r_I$ 至少为第一槽的行亏损；由 \eqref{hm:cross-power}
及 \eqref{hm:mode}，\eqref{hm:transport-identity} 给出
\[
 k_s^{-1}\sum_i(1-\rho_i^{m-2})\le C\|G_A-G_B\|^2+o(1).
\]
固定 $\rho_0\in(1/\sqrt2,1)$。由于 $2\rho_0^2>1+\eta_d$，绝对相关超过 $\rho_0$ 的边构成部分双射。
将其补为排列并逐行定向，使对角相关非负。未匹配行的比例由上一式除以
$1-\rho_0^{m-2}$ 控制；已匹配行使用 $1-x^m\le\frac m{m-2}(1-x^{m-2})$。
令 $\ell_i=1-\langle a_i,b_i\rangle^m$，遂得
\begin{equation}\label{hm:row-loss}
 k_s^{-1}\sum_i\ell_i+k_s^{-1}\sum_i\|a_i-b_i\|^2
 \le C\|G_A-G_B\|^2+o(1).
\end{equation}
定向时同步置换、变号学生系数以保持 $G_B$ 不变，模态界亦保持。
在此表示中 $\mathsf K_{II}=\prod_{a=1}^q(1-\ell_{i_a})/(1+e_d)^q\ge0$，
且 $1-\mathsf K_{II}\le\sum_a(\ell_{i_a}+e_d)$。
保留 \eqref{hm:transport-identity} 的对角平方项，记两系数的模态矩阵为 $H_C,H_{D'}$，得到
\[
 \begin{aligned}
 \|C-D'\|^2
 &\le\mathcal E+\sum_I(1-\mathsf K_{II})(c_I-d_I)^2\\
 &\le\mathcal E+2q\sum_i(\ell_i+e_d)
       ((H_C)_{ii}+(H_{D'})_{ii})
 \le C\|G_A-G_B\|^2+o(1).
 \end{aligned}
\]
最后去除一致趋于一的半径因子，并用 \eqref{hm:sym-group}，得到
\begin{equation}\label{hm:recover}
 \|C_s^A-C_s^B\|^2
 \le C\|G_A-G_B\|^2+o(1)
 \le C\|C_{s-1}^A-C_{s-1}^B\|^2+o(1).
\end{equation}

\subsection{内部符号传播与终层对齐}
\newthought*{终层只需匹配参数。}
内部层 $s<L$ 的有符号行置换写作 $S_s=D_sP_s$，其中 $D_s$ 对角元为 $\pm1$。隐特征相应变为
$R_sh^{(s)}$，$R_s=D_s^{m\bmod2}P_s$。同时以 $V^{(s+1)}R_s^{-1}$ 替换下一层矩阵，整个齐次函数精确保持不变。奇数 $m$ 将符号传播到下一层列，偶数 $m$ 仅传播排列；形式下游系数相应重排、变号，前述典型事件保持。

终层也允许有符号排列，但仅用作参数对齐标签。若匹配符号为 $\varsigma_i$、学生半径为 $t_i$，
把学生行定向后，同时在分析表示中取
$\overline c_i^B=\varsigma_i^{m}t_i^m/\sqrt{k_L}$，教师系数为
$\overline c_i^A=r_i^m/\sqrt{k_L}$。这样表示的原学生张量不变，模态对角界仍成立。
此时 $q_L=1$，由 \eqref{hm:sym-group} 和 \eqref{hm:row-loss}，定向后的终层平均行距离
直接由 $C\|C_{L-1}^A-C_{L-1}^B\|^2+o(1)$ 控制。无需恢复终层系数，
也不要求这个终层标签保持网络函数。

\subsection{深度归纳与任意先验耦合}
\newthought*{完成定理的证明}
令 $E_0=\|C_0(W)-C_0(V)\|^2$。自下而上在内部层选择有符号匹配、传播至下一层列，记当前对齐后的形式张量误差为 $E_s$。内部对称变换保持当前 $C_{s-1}$，且以后对齐不再改变已完成的较低层。因此 \eqref{hm:sym-group}、\eqref{hm:row-loss}、\eqref{hm:recover} 在每个内部层给出
\[
 E_s\le CE_{s-1}+e_d,\qquad
 \frac1{k_s}\sum_i\|w_{s,i}-(gv)_{s,i}\|^2\le CE_{s-1}+e_d,
 \qquad e_d\to0.
\]
半径误差由
$\|r_ia_i-t_ib_i\|^2\le C\|a_i-b_i\|^2+C\delta_d^2$
吸收。终层使用上一小节，固定深度归纳得到每层平均归一化行距离至多 $C_{L,m}E_0+o(1)$。
恢复原始权重需要乘 $k_{s-1}$，故以 $k_{s-1}k_s/p$ 加权求和，得到联合典型事件上的 \eqref{hm:main} 确定性版本。

具体映射族为
\[
 (gV)^{(1)}=S_1V^{(1)},\qquad (gV)^{(s)}=S_sV^{(s)}R_{s-1}^{-1}\quad(s\ge2),
\]
所有 $S_s$ 均可为有符号排列，终层标签不必是函数对称。因此
\[
 \log|\mathcal G_d|\le\sum_s\{\log(k_s!)+k_s\log2\}=O(d\log d).
\]
可固定标签次序打破所有平局，获得可测的匹配选择。任意先验边缘耦合下，两个各自典型事件的并集界保证联合失败概率趋零，不需要两个网络相互独立。坏事件上取恒等对齐，
\[
 p^{-1}\min_g\|W-gV\|^2\le2p^{-1}(\|W\|^2+\|V\|^2).
\]
归一化高斯权重平方范数有一致二阶矩；Cauchy--Schwarz 将坏事件期望变为 $o(1)$，且对耦合一致。取期望即证 \eqref{hm:main}。最后最高输入次数 $N=m^L$ 的高斯混沌满足
\[
 N!\|C_0(W)-C_0(V)\|^2
 =\|\Pi_N(\mathcal H_W-\mathcal H_V)\|_{L^2_X}^2
 \le\|\mathcal H_W-\mathcal H_V\|_{L^2_X}^2,
\]
证明函数距离版本。\hfill$\square$


\paragraph{实际函数的参数比较。}
定理 \ref{fp:global-function-pairing} 和正文的极化恒等式将实际函数距离控制转为齐次系数距离控制，再应用定理 \ref{hm:global-matching}。行符号只用于有限标签参数对齐，并不要求保持实际激活的函数值。

\section{条件 I--MMSE 与统一平滑测试场}\label{app:information}
\begin{lemma}[条件 I--MMSE]\label{info:immse}
设 $S\in\mathbb R^M$ 有有限二阶矩，$N\sim N(0,I_M)$ 独立于 $(S,Z)$。
对 $Y_s=\sqrt sS+N$，有
\[
 I(S;Y_\rho\mid Z)=\frac12\int_0^\rho
 \mathbb E\|S-\mathbb E[S\mid Z,Y_s]\|^2\,ds.
\]
\end{lemma}
\begin{proof}
条件于 $Z=z$，令 $q_v$ 为 $S+\sqrt vN$ 的密度。
Gaussian 热方程及分部积分给
\[
 \frac d{dv}h(q_v)=\frac12\mathcal J(q_v),\qquad
 \nabla\log q_v(y)=\frac{\mathbb E[S\mid z,S+\sqrt vN=y]-y}{v}.
\]
将 $S-(S+\sqrt vN)$ 对观察作正交分解，得到
$\operatorname{mmse}(S\mid z,S+\sqrt vN)=Mv-v^2\mathcal J(q_v)$。
从 $h(q_v)$ 减去噪声熵 $(M/2)\log(2\pi ev)$ 后微分，
再换元 $v=1/s$，即得互信息导数为 MMSE 的一半。
对有界信号可先在有界空间积分，再去掉空间截断；
一般信号用二阶矩截断逼近。Gaussian 卷积满足 $\mathcal J(q_v)\le M/v$，
截断信号的二阶矩一致可积，故每个紧区间 $0<s_0\le s\le s_1$ 上的积分恒等式通过极限。
最后用 $0\le I(S;Y_s\mid Z)\le s\mathbb E\|S\|^2/2$
及 $\operatorname{mmse}\le\mathbb E\|S\|^2$，平均 $z$ 并令 $s_0\downarrow0$。
\end{proof}

\subsection{真实通道的分部积分界}
记 $f=F_{\theta(W)}^{(L)}(X)$、$a=\nabla_Wf$、
$\mathcal I_W=\E_Xaa^T/\Delta$、$Z=(\mathcal D_n,X)$，
$U_s=\sqrt sf+N$，其中 $N$ 独立标准 Gaussian。上引理给
\begin{equation}\label{eq:conditional-immse}
 \epsilon_{\log,d,n}=\frac12\int_0^{1/\Delta}m(s)\,ds,
 \qquad m(s)=\E(f-\E[f\mid Z,U_s])^2.
\end{equation}
对只依赖 $(W,X)$ 的向量场 $b$，令
\[
 G_s=\delta b+b^TS_D+\sqrt sN b^Ta,\qquad
 \delta b=\operatorname{div}_Wb-W\cdot b,\quad
 S_D=\nabla_W\log p(\mathcal D_n\mid W).
\]
固定观察并对参数分部积分得
$\E[(\E[f\mid Z,U_s]-f)G_s]=\E b^Ta$。
给定 $(W,X)$，训练得分与辅助噪声独立且中心化，训练得分协方差为
$n\mathcal I_W$。因此 Cauchy--Schwarz 给
\begin{equation}\label{eq:vt-mmse}
 m(s)\ge\frac{(\E b^Ta)^2}
 {n\E b^T\mathcal I_Wb+\E(\delta b)^2+s\E(b^Ta)^2}.
\end{equation}
先对紧支撑光滑场及截断的后验均值证明，再用引理 \ref{an:sobolev}
的 Sobolev 逼近和 $L^2$ 收敛取极限；以下估计保证所需可积性。

\subsection{两支共用的正则化场及其导数}
固定 $0<r,\dsm\le1$，令 $a_{\dsm}=\nabla_WF_{\dsm}(W,X)$、
$\mathcal I_{\dsm}=\E_Xa_{\dsm}a_{\dsm}^T/\Delta$、
$\lambda=r/p$、$M=(\mathcal I_{\dsm}+\lambda I)^{-1}$、$b=Ma_{\dsm}$。
仅测试场被平滑，$f$、$U_s$ 和训练似然均仍为真实网络。
定义算子 $T,T_{\dsm}:\mathbb R^p\to L^2(P_X)$：
$(Tu)(X)=a(X)^Tu/\sqrt\Delta$，$T_{\dsm}$ 同理。
于是 $T^*T=\mathcal I_W$、$T_{\dsm}^*T_{\dsm}=\mathcal I_{\dsm}$。

为控制逆矩阵导数，记 $C_i=M^{1/2}(\partial_i\mathcal I_{\dsm})M^{1/2}$。
由 $\|T_{\dsm}M^{1/2}\|_{\rm op}\le1$，有
\[
 C_i=(\partial_iT_{\dsm}M^{1/2})^*T_{\dsm}M^{1/2}
       +(T_{\dsm}M^{1/2})^*\partial_iT_{\dsm}M^{1/2},\qquad
 \|C_i\|_{\rm HS}^2\le4\|M\|_{\rm op}\|\partial_iT_{\dsm}\|_{\rm HS}^2.
\]
保留 $\partial_iM=-M(\partial_i\mathcal I_{\dsm})M$，并用
\[
 \E_X\|M(\partial_i\mathcal I_{\dsm})Ma_{\dsm}\|^2
 \le\Delta\|M\|_{\rm op}\|C_i\|_{\rm HS}^2,
\]
对 $i$ 求和，得到
\[
 \E\|Db\|_{\rm HS}^2
 \le C\lambda^{-2}\E\|\nabla_W^2F_{\dsm}\|_{\rm HS}^2
 \le Cp^2/(r^2\dsm).
\]
又 $\E\|b\|^2=\Delta\E\tr(M^2\mathcal I_{\dsm})\le\Delta p/(4\lambda)$。
引理 \ref{an:sobolev} 及梯度四阶矩遂给
\begin{equation}\label{unified:field-cost}
 \E(\delta b)^2\le Cp^2/(r^2\dsm),\qquad
 \E(b^Ta)^2\le\lambda^{-2}\E(\|a_{\dsm}\|^2\|a\|^2)\le Cp^2/r^2.
\end{equation}

\subsection{比较真实与平滑信息矩阵}
记 $E=T-T_{\dsm}$，则 \eqref{eq:smoothing} 给
$\E\|E\|_{\rm HS}^2\le C\dsm$。令
$M_0=(\mathcal I_W+\lambda I)^{-1}$、
$D_{d,\dsm}(r)=(\lambda/p)\E\tr M$。
因为
\[
 M-M_0=M(T_{\dsm}^*E+E^*T)M_0,
\quad
 \|T_{\dsm}M\|_{\rm HS},\ \|TM_0\|_{\rm HS}
 \le\sqrt{p/(4\lambda)},\quad
 \|M\|_{\rm op},\ \|M_0\|_{\rm op}\le1/\lambda,
\]
对两个乘积逐项取迹并用 Hilbert--Schmidt Cauchy--Schwarz，得到
\[
 |D_{d,\dsm}(r)-D_d(r)|
 \le\frac{\E\|E\|_{\rm HS}}{\sqrt{p\lambda}}
 \le C\sqrt{\dsm/r}.
\]
同样由 $\E b^Ta=\Delta\E\tr(MT_{\dsm}^*T)$，有
\[
 \E b^Ta\ge\Delta p[1-D_{d,\dsm}(r)]
       -\Delta\sqrt{p/(4\lambda)}(\E\|E\|_{\rm HS}^2)^{1/2}
 \ge\Delta p[1-D_d(r)-C\sqrt{\dsm/r}].
\]
最后在 $L^2(W)$ 中使用算子 Hilbert--Schmidt 范数的三角不等式。
由 $\|T_{\dsm}MT_{\dsm}^*\|_{\rm HS}^2\le p$ 及
$\|EMT_{\dsm}^*\|_{\rm HS}^2\le\|E\|_{\rm HS}^2/(4\lambda)$，得到
\[
 (\E b^T\mathcal I_Wb/\Delta)^{1/2}
 =(\E\|TMT_{\dsm}^*\|_{\rm HS}^2)^{1/2}
 \le\sqrt p+C\sqrt{\dsm/\lambda}
 =\sqrt p(1+C\sqrt{\dsm/r}).
\]
将目标取正部，把上述估计及 \eqref{unified:field-cost} 代入
\eqref{eq:vt-mmse}，再对 \eqref{eq:conditional-immse} 积分，即得
\eqref{eq:lower-unified}。精确系数 $1/2$ 来自真实通道的条件 I--MMSE，
而不是后验 Gaussian 近似。

\end{fullwidth}
\end{document}
